\documentclass{article}

% if you need to pass options to natbib, use, e.g.:
    \PassOptionsToPackage{numbers, compress}{natbib}
% before loading neurips_data_2023

% ready for submission
% \usepackage{neurips_data_2023}

% to compile a preprint version, add the [preprint] option, e.g.:
    \usepackage[final]{neurips_data_2023}
% This will indicate that the work is currently under review.

% to compile a camera-ready version, add the [final] option, e.g.:
%     \usepackage[final]{neurips_data_2023}

% to avoid loading the natbib package, add option nonatbib:
%    \usepackage[nonatbib]{neurips_data_2023}

% Submissions to the datasets and benchmarks are typically non anonymous,
% but anonymous submissions are allowed. If you feel that you must submit 
% anonymously, you can compile an anonymous version by adding the [anonymous] 
% option, e.g.:
    % \usepackage[anonymous]{neurips_data_2023}
% This will hide all author names.

\usepackage[utf8]{inputenc} % allow utf-8 input
\usepackage[T1]{fontenc}    % use 8-bit T1 fonts
\usepackage{hyperref}       % hyperlinks\
\usepackage{url}            % simple URL typesetting
\usepackage{booktabs}       % professional-quality tables
\usepackage{amsfonts}       % blackboard math symbols
\usepackage{nicefrac}       % compact symbols for 1/2, etc.
\usepackage{microtype}      % microtypography
\usepackage[dvipsnames]{xcolor} % colors
\usepackage{enumitem}
\usepackage{graphicx}
\usepackage{tabularx}
\usepackage{colortbl}
\usepackage{dirtree}
\usepackage{subcaption}
\usepackage{amsmath}
\usepackage{float}
\usepackage{listings}
\usepackage{calc}
% Custom Colors List
\definecolor{light-gray}{gray}{0.95}

% Custom Commands List
\newcommand{\benchmark}{InterCode}
\newcommand{\code}[1]{\colorbox{light-gray}{\texttt{#1}}}
\newlength\myheight
\newlength\mydepth
\settototalheight\myheight{Xygp}
\settodepth\mydepth{Xygp}
\setlength\fboxsep{0pt}
\newcommand*\inlinegraphics[1]{%
  \settototalheight\myheight{Xygp}%
  \settodepth\mydepth{Xygp}%
  \raisebox{-\mydepth}{\includegraphics[height=1.5\myheight]{#1}}%
}

% Editing Commands
\newcommand{\todo}[1]{\textcolor{red}{[TODO: #1]}} 
\newcommand{\sy}[1]{\textcolor{blue}{[Shunyu: #1]}}
\newcommand{\ap}[1]{\textcolor{purple}{[Akshara: #1]}}
\newcommand{\kn}[1]{\textcolor{magenta}{[Karthik: #1]}}
\newcommand{\jy}[1]{\textcolor{orange}{[John: #1]}}

\title{\benchmark: Standardizing and Benchmarking Interactive Coding with Execution Feedback}

% The \author macro works with any number of authors. There are two commands
% used to separate the names and addresses of multiple authors: \And and \AND.
%
% Using \And between authors leaves it to LaTeX to determine where to break the
% lines. Using \AND forces a line break at that point. So, if LaTeX puts 3 of 4
% authors names on the first line, and the last on the second line, try using
% \AND instead of \And before the third author name.

\author{%
  John Yang \quad Akshara Prabhakar \quad Karthik Narasimhan \quad Shunyu Yao\\
%   Department of Computer Science\\
  Department of Computer Science, Princeton University\\
%   Pittsburgh, PA 15213 \\
  \texttt{\{jy1682, ap5697, karthikn, shunyuy\}@princeton.edu}
}

\begin{document}

\maketitle

\renewcommand{\thefootnote}{\fnsymbol{footnote}}


\begin{abstract}
Humans write code in a fundamentally interactive manner and rely on constant execution feedback to correct errors, resolve ambiguities, and decompose tasks. While LLMs have recently exhibited promising coding capabilities, current coding benchmarks mostly consider a static instruction-to-code sequence transduction process, which has the potential for error propagation and a disconnect between the generated code and its final execution environment. To address this gap, we introduce \benchmark{}, a lightweight, flexible, and easy-to-use framework of interactive coding as a standard reinforcement learning (RL) environment, with code as actions and execution feedback as observations. Our framework is language and platform agnostic, uses self-contained Docker environments to provide safe and reproducible execution, and is compatible out-of-the-box with traditional seq2seq coding methods, while enabling the development of new methods for interactive code generation.
% \benchmark{}'s execution-oriented evaluation offers a multitude of novel feedback signals that can be used to develop interactive coding agents that leverage the extra contextual information.
% \benchmark{} defines a general formulation for \textit{interactive coding tasks} and provides a framework that makes it easy to construct such a task setting from existing, static instruction-to-code datasets.
We use \benchmark{} to create three interactive code environments with Bash, SQL, and Python as action spaces, leveraging data from the static NL2Bash~\cite{lin-etal-2018-nl2bash}, Spider~\cite{yu-etal-2018-spider}, and MBPP~\cite{austin2021program} datasets. We demonstrate \benchmark{}'s viability as a testbed by evaluating multiple state-of-the-art LLMs configured with different prompting strategies such as ReAct~\cite{yao2023react} and Plan \& Solve~\cite{wang2023planandsolve}. Our results showcase the benefits of interactive code generation and demonstrate that \benchmark{} can serve as a challenging benchmark for advancing code understanding and generation capabilities. \benchmark{} is designed to be easily extensible and can even be used to create new tasks such as Capture the Flag, a popular coding puzzle that is inherently multi-step and involves multiple programming languages.
\footnote{Code and data available at \url{https://intercode-benchmark.github.io/}}
% % Finally, we develop a new Capture the Flag task with \benchmark{} that showcases \benchmark{} value towards constructing new tasks, datasets, and evaluation metrics.
% \ap{
% that showcases the potential of \benchmark{} in enabling the development of innovative tasks and metrics.}
\end{abstract}
\section{Introduction}
\label{sec:intro}

% Problem with existing landscape of code + NL benchmarks
% \begin{itemize}[noitemsep,topsep=0pt]
%     \item Code Benchmarks like MBPP, CodeRL provide full specifications up front - this covers algorithm-solving style challenges, but practical code tasks are very different + require much more context.
%     \item New tasks, prompting schemes tend to evaluate specific reasoning capabilities, with a task well-suited for specific prompting strategies
% \end{itemize}

% Two aspects of motivation:
% \begin{itemize}[noitemsep,topsep=0pt]
%     \item Solve coding problems. Existing benchmarks are mostly (if not all) static. While some models leverage  execution~\cite{alphacode} / iterative-feedback~\cite{shinn2023reflexion}/ decomposition~\cite{parscal}, each method has its own setup of using interaction/execution. need a benchmark 
%     \item (less important motivation) Interesting interaction. Code is an interesting action space (compositional, diverse, semantic, ...), different from all existing action types in various RL/interactive environments: low-level control, natural language, etc. there has been 
% \end{itemize}

The art of computer programming is naturally an interactive process. When a human programmer writes code, she relies on several iterations of a `write-execute-test' loop in order to iteratively refine solutions, plan changes, test sub-modules, and solve ambiguities by checking execution behavior. While this is reminiscent of other human endeavors like writing, code compilation and execution produce exact results that provide a deterministic form of feedback to make the refinement process more straightforward. Depending on the observed results, programmers perform various levels of debugging and rewriting, and continue the process until their code satisfies the requirements.

There has been increasing interest in recent years around the development of models that can automatically generate code given a specification in natural language~\cite{feng2020codebert, wang2021codet5, clement2020pymt5, alphacode, le2022coderl}. Powered by large-scale pre-training over thousands of codebases~\cite{Agashe2019JuICeAL, husain2020codesearchnet, gao2020pile}, these models have shown solid performance on static benchmarks like HumanEval~\cite{chen2021evaluating}, APPS \cite{hendrycksapps2021}, MBPP \cite{austin2021program}, CodeXGLUE \cite{lu2021codexglue}. However, generating code in a static, sequence-to-sequence or auto-regressive fashion has several drawbacks: 1) simple errors (even typos) can propagate and there is no chance for recovery or revision, 2) there is a disconnect between the code generation process and its downstream execution on the desired software and hardware environment, and 3) there is little room for human intervention or collaboration in the code generation process. 

\begin{figure}
    \centering
    \includegraphics[width=1.0\textwidth]{figures/assets/Preview.pdf}
    \caption{Overview of \benchmark. Setting up an interactive code environment with \benchmark{} requires a Dockerfile, dataset, reward function definition, and a small amount of subclass implementation. The interactive loop between agent and environment closely mirrors real world software development processes. While \benchmark{} task performance is generally quantified as a binary 0/1 completion score, \benchmark{} allows for the design of more complex evaluation criteria that can incorporate execution output and the effects of interaction on the state space.}
    \label{fig:preview}
    \vspace{-1.5em}
\end{figure}

Recently, some works have proposed the use of execution feedback or interaction~\cite{wang2023interactive} to benefit code generation models~\cite{Lai2022DS1000, huang-etal-2022-execution, wang2023executionbased, hendrycksapps2021}. However, these papers consider their own individual setup and are difficult to compare with one other due to the use of different compilers, execution environments, feedback signals, and assumptions on the interactive process such as human participation to create task descriptions or provide natural language feedback. This makes it difficult to compare existing methods for code generation and to clearly understand the benefits of interactive generation. 

To address these issues, we propose InterCode, the first standard coding benchmark designed natively with an interactive execution environment. Closely mimicking the human decision-making process, \benchmark{} allows a coding agent to interactively receive feedback from compilers/interpreters that execute its code, and to submit further refinements. We design \benchmark{} to be like a standard reinforcement learning (RL) environment that requires minimal human intervention and one in which generated code is treated as actions, which are executed to reveal observations. Our framework is (1) language and platform agnostic and can easily be used for new coding problems, (2) uses self-contained Docker environments to provide safe execution, and (3) compatible out-of-the-box with traditional seq2seq generation methods, while also enabling and empowering the development of new interactive techniques. 

We demonstrate the power of the framework by implementing Bash, SQL, and Python tasks within \benchmark{}, building on pre-existing static datasets~\cite{zhong2017seq2sql, lin-etal-2018-nl2bash, austin2021program}. We perform experiments across diverse models and prompting methods, including ReAct~\cite{yao2023react} and Plan \& Solve~\cite{wang2023planandsolve}. Our findings concretely showcase the benefits of interaction towards solving coding tasks, discuss the distribution of distinct code understanding challenges across different task settings, and explore the ease with which new tasks and datasets can be defined using \benchmark{}.

To summarize, our paper makes the following contributions:
\begin{itemize}[noitemsep,topsep=0pt]
\item We develop \benchmark{}, a new, universal framework for interactive code generation, which provides ease of use, extensibility, and safety.
\item Using \benchmark{}, we perform a comprehensive evaluation of state-of-the-art models and identify several avenues for improvements.
\item We release our framework as a new benchmark along with useful empirical tools to customize any new static code datasets into interactive tasks.
\end{itemize}

% \sy{raw paragraph-by-paragraph plan}

% Human coding is interactive by nature. Execution is the most nature interaction people leverage, for all sorts of uses: iterative debug, decompose and plan, figure out uncertainty, internally evaluate, etc...

% Traditionally, coding benchmarks have been seq2seq. But recently, there has been numerous methods leveraging execution feedback and/or interaction, e.g. xxx. However, each method has its own setup and hard to (fairly) compare, e.g. use different compilers or execution environments, different assumptions, ... And thus hard to understand differnet aspects of benefits interaction can bring.

% We propose InterCode, the first standard coding benchmark designed with an interactive execution environment. Formulated as a standard RL environment, where a coding agent interact with a coding execution environment. We implemented SQL and bash environments, and \sy{minimal description of what you actually do to build the benchmark, e.g. data -> env -> evaluation -> ...} 


% InterCode is designed to be \sy{adjust} (1) standard RL env with minimal human intervention and "code-as-action" (2) easy to plug in new coding problems (3) safe (docker) (4) compatiable with traditional seq2seq coding methods, while enabling brand new interactive methods... 







\section{Related Work}
\label{sec:related}

\textbf{Interactive environments for coding.} 
Most coding benchmarks (e.g.\,SQL - Spider \cite{yu-etal-2018-spider}, KaggleDBQA \cite{lee-etal-2021-kaggledbqa}; Bash - NLC2CMD \cite{pmlr-v133-agarwal21b}, NL2Bash \cite{lin-etal-2018-nl2bash}; Python - HumanEval \cite{chen2021evaluating}, APPS \cite{hendrycksapps2021}, MBPP \cite{austin2021program}, CodeXGLUE \cite{lu2021codexglue}, CodeNet \cite{puri2021codenet}) frame the coding problem as a sequence transduction problem (from instruction to code), rather than an interactive decision making problem with an execution environment. 
Attempts have been made to simulate interaction by developing conversational, dialogue-style \cite{yu-etal-2019-sparc, yu2019cosql}, multi-step problem solving \cite{nijkamp2022codegen} datasets, which involve pre-annotated human-designed queries. The work closest to \benchmark{} has been recent explorations of Python Jupyter Notebooks as a natural choice for interactive coding \cite{huang-etal-2022-execution, Lai2022DS1000, yin2022natural}. However, task data and settings often constrain allowed actions to a closed domain of code and libraries~\cite{Lai2022DS1000, yin2022natural}, use evaluation procedures or metrics that may not generalize~\cite{huang-etal-2022-execution}, require human-in-the-loop participation (i.e. create task contexts, write problems, evaluate execution per task instance)~\cite{Lai2022DS1000}, or are Python-exclusive~\cite{huang-etal-2022-execution, Lai2022DS1000, yin2022natural, wang2023executionbased}.
% In addition, such benchmarks cannot be adapted to other programming languages or task settings in a straightforward manner.
\benchmark{} provides a more general purpose foundation defining interactive coding tasks that enables easy construction of diverse task settings, can have any programming language(s) as the action space, and has automatic, execution-based evaluation.

\textbf{Execution-based evaluation for coding.} Evaluation for NL-to-code generation models has recently shifted away from surface form similarity metrics (BLEU~\cite{Papineni2002BleuAM, Agashe2019JuICeAL}, ROUGE~\cite{lin-2004-rouge}, Exact Match) towards execution oriented ratings (unit tests~\cite{austin2021program, chen2021evaluating, huang-etal-2022-execution, Lai2022DS1000, hendrycksapps2021}, output matching~\cite{dong2016language, huang-etal-2022-execution, zhong2017seq2sql}).
The rigidity of surface form analysis overlooks code syntax features, ignores execution effect, or over-penalizes alternative solutions~\cite{zhou2023codebertscore},
On the contrary, execution-based assessment is a more thorough and comprehensive score of code functionality~\cite{hendrycksapps2021} and is a more natural fit for open-domain program usage that does not constrain code generation to a subset of the language space~\cite{wang2023executionbased}.
However, for newer benchmarks and datasets that put forth task definitions incorporating execution-based evaluation (APPS~\cite{hendrycksapps2021}, ExeDS~\cite{huang-etal-2022-execution},  ODEX~\cite{wang2023executionbased}), the fundamental code generation task (Context + Code $\rightarrow$ Execution $\rightarrow$ Score) is still devoid of interaction.
% In addition, the proposal of new metrics borne from different task formulations may make it difficult to compare models across all settings.
\benchmark{} combines execution-based evaluation with flexible task construction, enabling more diverse problem-solving paradigms within a unified coding task formulation.
\benchmark{}'s use of virtual containers as execution sandboxes protect against harmful actions and allow for advanced evaluation criteria beyond the aforementioned ones.
% surface form, unit tests, or execution output matching.
% An example is the \benchmark{}-Bash environment, which includes file system modification in its task completion function, a capability not easily achievable with previous evaluation methods.

\textbf{Methods for interactive or execution-based coding.} The value of generative code models and interactive problem solving has motivated a recent proliferation of work to augment reasoning capabilities' of existing language models~\cite{yao2023react, shinn2023reflexion, wang2023planandsolve, yao2023tree, zhang2023planning, chen2023teaching} or propose new modeling techniques to tackle coding as a sequential decision making and reasoning tasks~\cite{bunel2018leveraging, Chen2018ExecutionGuidedNP, ellis2019write, alphacode, chen2022codet, le2022coderl}, of which evaluation is unit test based.
Approaches that leverage execution typically use re-ranking~\cite{zhang2022coder, ni2023lever, yin-neubig-2019-reranking, zeng2022nbest} or majority vote~\cite{Chen2018ExecutionGuidedNP, alphacode, shi2022natural} to decide on a final prediction.
Additional work also explores incorporating human-in-the-loop~\cite{chen2023improving, lahiri2022interactive}, compilers~\cite{wang2022compilable}, and text~\cite{wang2023leti,zhang2023selfedit} feedback.
A common thread among these contributions is that 1) the task setting can only provide the investigated form of feedback and 2) sought-after capabilities are exemplified by strong performance on favorably curated tasks and datasets, rendering comparisons across benchmarks tedious.
\benchmark{} has the potential to standardize the evaluation of these methods because 1) the interactive coding task is a conglomeration of many interesting interaction, reasoning, and decision-making challenges and 2) \benchmark{}'s task construction makes it possible to incorporate a wide variety of sources of feedback.

% \textbf{Other language-based interaction} Text games, web navigation, etc. but natural langauge, not programming language as action. \sy{less important, hold for now} 

% \textbf{Interactive coding benchmarks}
% % existing benchmarks
% %As models become increasingly efficient, developing more challenging and general code generation benchmarks has become the norm. 
% Most coding benchmarks (e.g.\,SQL \cite{yu-etal-2018-spider,lee-etal-2021-kaggledbqa}, bash \cite{pmlr-v133-agarwal21b, lin-etal-2018-nl2bash}, Python \sy{explicitly write down names}) set up a coding problem as an instruction and a .

% In the absence of an interactive coding benchmark that can cater to more than just Python, we develop \benchmark{}, which currently provides interactive environments for SQL, Bash, and Python and demonstrate that it can be extended to other exciting code-based tasks like CTFs.



% \textbf{Interactive or execution-based coding methods} Several past and recent works employ variations of test-driven development
% or code debugging practices. AlphaCode~\citep{li2022competition}
% evaluates a set of generations on hidden test cases. CodeT~\citep{chen2022codet}
% uses self-generated unit tests that are used to score generated function
% implementations.
% Self-Debugging~\citep{chen2023teaching} employs a debugging component that is used
% to improve existing implementations given feedback from a code
% execution environment. CodeRL~\citep{le2022coderl} sets the problem
% in an RL framework using an actor-critic setup to debug programs given
% feedback from an execution environment. AlphaCode, Self-Debugging and CodeRL
% are effective in fixing less-complex program bugs, but they rely
% upon ground truth test cases that invalidate pass@1 eligibility,
% and do not use self-reflection to bridge the gap between error
% identification and implementation improvement. CodeT does not access
% hidden test cases but does not implement a self-learning step to
% improve code writing.




% The standard approach to generating code involves models generating a blob of code given a natural language intent and are typically tested via hand-designed test cases \cite{chen2021evaluating, alphacode,hendrycksapps2021}.
% % multi-turn: decompose
% Progressive generation of subprograms through multi-turn expert-generated prompts \cite{nijkamp2022codegen} demonstrated that Language Models could benefit from the systematic multi-step approach to code generation but offers no scope for error-recovery in case of incorrectly generated previous turns. 
% Recently, language models have emerged as agents capable of observing, interacting, and receiving iterative feedback from external entities and the environment \cite{wang2023interactive}. 

% - Planning \cite{zhang2023planning}

% - Human-in-the-loop \cite{chen2023improving, lahiri2022interactive}

% - Compiler feedback \cite{wang2022compilable}

% - Textual feedback \cite{wang2023leti}



% \textbf{Language-based Interaction} \sy{less important} 
% % \textbf{Text Games}
% % CALM: Similar challenge of determining both admissible and promising actions when the action space is extremely large (essentially infinite)

\section{The \benchmark{} Benchmark}
\label{sec:benchmark}

% We present the \benchmark{} benchmark, a framework for designing interactive code environments. We first discuss the task formulation put forth by \benchmark{}, then dive into its implementation and how it supports versatile tasks. Finally, we demonstrate \benchmark{}'s potential to bring forth the variety of language grounding and understanding challenges that exist in the space of interactive coding.

\subsection{Formulation}
\label{sec:benchmark:formulation}

The \benchmark{} benchmark formalizes interactive coding with execution feedback as a partially observable Markov decision process (POMDP) $(\mathcal{U}, \mathcal{S}, \mathcal{A}, \mathcal{O}, \mathcal{T}, \mathcal{R})$ with instruction space $\mathcal{U}$, state space $\mathcal{S}$, action space $\mathcal{A}$, observation space $\mathcal{O}$, transition function $\mathcal{T}: \mathcal{S} \times \mathcal{A} \to \mathcal{S}$, and reward function $\mathcal{R}: \mathcal{S} \times \mathcal{A} \to [0, 1]$.
% Mirroring real-world coding paradigms, the primary interactive setup presented by \benchmark{} is, 
Given a  coding instruction $u \in \mathcal{U}$ in natural language, an agent issues code or a special \texttt{submit} keyword as an action $a_t \in \mathcal{A}$.
An action is \textit{admissible}~\cite{yao2020calm} if it can be parsed and executed in the compiler/interpreter environment, and an admissible action incurs a change in the latent state space $s_{t+1} \in \mathcal{S}$, and an execution feedback as observation $o_{t+1} \in \mathcal{O}$.
The interaction loop repeats until the \texttt{submit} action is issued, wherein the task episode ends and a reward $r = \mathcal{R}(s_T, \texttt{submit}) \in [0, 1]$ is computed, with $1$ representing task completion. 
We use the \textbf{Success Rate (SR)} metric, defined as the proportion of task episodes where $r = 1$.
We also define the \textbf{Error \%} metric, which is the percentage of \textit{non} admissible actions across  task episodes.

% A state $s_t \in \mathcal{S}$ generally represents the system or setting that the interactive coding task is taking place in (i.e. MySQL Database, Bourne Shell). Depending on the system, the state includes information such files/tables/folders, permissions, current working directory, execution context, and details about the OS or installed dependencies, to name a few. This latent state is not directly available to an agent, but can be uncovered via observations from agent actions. 


\subsection{Construction pipeline}
\label{sec:benchmark:construct_pipeline}
At a high level, \benchmark{} decomposes the construction of an interactive coding task into three \textbf{modular} parts: (1) environment construction, (2) data collection, and (3) reward design.
This workflow allows for the safe execution of transition functions, flexible reward design, and convenient adaptation of existing instructions to an interactive setting.

\textbf{Docker-based environments.}
\benchmark{} uses Docker~\cite{merkel2014docker} virtual containers as a general-purpose execution sandbox.
Given a \texttt{Dockerfile} that defines a system and execution entrypoint, \benchmark{} creates a corresponding, stateful virtual container that hosts the desired state space and transition function.
We choose Docker as the basis of \benchmark{}'s environment construction for its safe execution in virtual containers, reproducibility of a \texttt{Dockerfile} across any Docker-equipped machine, and excellent coverage of application code, libraries, and dependencies offered by the \texttt{Dockerfile} DSL.

\textbf{Data collection.}
\benchmark{} requires that a dataset has at minimum two fields: \texttt{query}, a natural language instruction $u \in \mathcal{U}$, and \texttt{gold}, an answer or code block that is a procedure for generating the correct answer.
We define these conditions to make it easy to adapt existing text-to-code datasets to an interactive setting while also leaving plenty of bandwidth for constructing new tasks and datasets.

\textbf{Reward design.}
Across a single task episode, the action, observation, and state modification (if any) per interaction loop are implicitly logged by \benchmark{}.
\benchmark{}'s default reward function determines task completion via an exact match of the agent's execution output (observation and state modifications) against the gold command, where $1$ is awarded only if all components match.
Since Exact Match is usually too stringent of an evaluation criteria, \benchmark{} exposes a reward function endpoint that has access to both the interaction history and the execution container, allowing for custom reward function definitions that can incorporate multiple signals.

\subsection{Implementations}
\label{sec:benchmark:ic_envs}

Following the procedure discussed in Section~\ref{sec:benchmark:construct_pipeline}, we create two separate \benchmark{} based environments where Bash and SQL are the action spaces respectively. Table~\ref{table:env-compare} summarizes them.

\begin{table}
\centering
\begin{tabularx}{\textwidth}{lllll}
\toprule
    Action Space
    & Environment \inlinegraphics{figures/assets/emojis/docker.pdf}
    & Dataset \inlinegraphics{figures/assets/emojis/dataset.pdf}
    & Reward Function \inlinegraphics{figures/assets/emojis/reward.pdf} \\
\midrule
    Bash & Ubuntu Terminal & NL2Bash~\cite{lin-etal-2018-nl2bash} (200) & Latest Std. Output + File System $\Delta$ \\
    SQL & MySQL Database & Spider 1.0~\cite{yu-etal-2018-spider} (1034) & Latest Std. Output \\
    Python & Python Interpreter & MBPP~\cite{austin2021program} (117)   & Submitted Function \\
\bottomrule
\end{tabularx}
\vspace{0.5em}
\caption{Rundown of the two environments with Bash and SQL as action spaces developed using the \benchmark{} framework. The numbers in parentheses refer to the number of task instances adopted from each dataset. Each environment is defined in under 200 lines of code total. Specific discussion of the environment construction and reward function can be found in \S~\ref{appx:env:bashenv} and \S~\ref{appx:env:sqlenv}}
\label{table:env-compare}
\vspace{-1.5em}
\end{table}

\textbf{\benchmark{}-Bash.} We define a \texttt{bash} shell within an Ubuntu Operating System as the task setting.
To evaluate an agent's ability to adapt generations to different situations, we architect four distinct file systems that can be swapped into the Bash environment by changing a single line in the \texttt{Dockerfile}.

We bootstrap the NL2Bash~\cite{lin-etal-2018-nl2bash} dataset (which lacks specificity in queries and grounding to any underlying file system, preventing it from being used directly for interactive evaluations) to create an interactive coding task where an agent completes an instruction via \texttt{bash} actions. Transferring NL2Bash to the interactive task setting requires simple transformations to ground instructions and \texttt{gold} code blocks in the file system. First, we consider a subset of 1000 commands with each having $\geq$ 4 utilities. We then filter out commands that are non-UNIX, non-Linux, or use utilities we currently do not support (eg. "ssh", "sudo", time, and GUI-dependent utilities). Finally, we enhance under-specified commands with specific file names/directory names/paths and update deprecated utilities/flags. The resulting 200 commands are grouped into 4 disjoint sets, 3 of which were grounded to custom-designed file systems, while one set is file-system agnostic. This categorization allows for a comprehensive evaluation of different command-grounding scenarios.

The \benchmark{}-Bash dataset instructions typically make one or both of the following two types of requests. It either
1. Requests information that can be answered via execution output (i.e. "How many files...", "What is the size of...", "Where is \texttt{<file>} stored?") or
2. Requests a change to the location/configuration/content of a file or folder (i.e. "Move \texttt{dir1} folder...", "Set permissions of...", "Append a line to...").
Therefore, we define a custom reward function that evaluates an agent's performance against file system modifications and the latest execution output.
Execution output is graded with a simple lexical similarity function.
File system assessment is done in two parts. First, a comparison of the agent's and \texttt{gold} command's list of file system changes (list of \texttt{[path, modification type $\in$ [added, changed, deleted]]} entries) reveals any extraneous or missing changes.
Second, \texttt{md5sum} hashes of each commonly edited file path are compared to determine if an added or changed file was altered correctly.
A max score of 1 is achieved only if the correct file paths are changed, the changes are correct, and the latest execution output matches the gold command output exactly. Additional Bash statistics and design details are discussed in \S~\ref{appx:env:bashenv}.

\textbf{\benchmark{}-SQL.} We write a \texttt{Dockerfile} that defines a SQL interpreter within a MySQL database as the task setting.
To create the databases and tables necessary for the task dataset, we write type resolution scripts and perform database conversions using the \texttt{sqlite3mysql}~\cite{sqlite3mysql} Python library to adapt the Spider~\cite{yu-etal-2018-spider} database and table schema to a MySQL format.
We then consolidate all setup code into a single, unified MySQL \texttt{.sql} dump that contains the complete set of schemas for all tables across 20 different databases.
On container start-up, this file is invoked automatically, creating and populating databases with tables and tables with records.

The Spider~\cite{yu-etal-2018-spider} dataset is a large-scale cross-domain dataset originally meant for evaluating SQL query generations from natural language questions. We adapt the development set, which contains $1034$ task instances, and remove all extraneous columns aside from the natural language questions and gold SQL command. The \texttt{instruction} and \texttt{gold} values do not require any additional pre-processing to be compatible with the MySQL task environment.

Finally, we employ Intersection over Union (\textit{IoU}), or more formally the Jaccard Index, to quantify the correctness of the latest execution output generated by the agent against the gold output, where both outputs are a list of records.
A non-tabular execution output receives a reward of 0 by default.
Among the items that lie in the intersection of the agent and gold execution outputs, we also apply a penalty if the records are in the incorrect order.
To quantify how sorted the agent output is relative to the gold output, we lean on Kendall’s $\tau$ and adjust the output range to $[0, 1]$. The \textit{IoU} score is then directly scaled by this coefficient. All in all, only a correctly ordered list with the exact set of records found in the gold output receives a score of 1.
Visualizations like Figure~\ref{fig:preview} for SQL along with a more extensive implementation discussion for this environment are in \S~\ref{appx:env:sqlenv}

\textbf{\benchmark{}-Python.} In this setting, we define a Python interpreter running within an Ubuntu operating System as the task setting.
The Dockerfile can be configured to run any Python version.
The interpreter is not initialized with any dependencies, but PyPI packages can be installed and used by the agent.

We use the MBPP~\cite{austin2021program} dataset which presents the code completion task of synthesizing Python code from a method header and docstring.
Evaluation of correctness is performed with an associated set of unit tests given by MBPP.
The MBPP dataset is straightforward to adapt to the interactive setting, requiring no modifications to the query or evaluation components.
Finally, we directly inherit MBPP's evaluation procedure of proportion of unit tests passed.
With \benchmark{}, it is easy to use existing datasets to evaluate how well models can use different programming languages as actions.

\textbf{Validations.}
To verify the functionality of action execution in the task environment and the correctness of custom reward functions, we write testing scripts for both Bash and SQL that pass the gold command in as a dummy agent's action to ensure that the command is admissible and executes without error, and to verify that the reward received by the command is $1$.
To confirm that \benchmark{}'s dataset specification is enforced across multiple accepted file formats, we define a custom \benchmark{} data loader class which is then rigorously unit tested.

% \subsection{Research Challenges Imposed by \benchmark{}}
% \label{sec:benchmark:challenges}
% \benchmark{}'s interface and design provides fertile ground for constructing language understanding and decision-making challenges spanning across NLP and RL.
% Solving a task across multiple iterations implicitly relies on a model's ability to perform \textit{error correction} by determining what action to perform next based on past execution output(s).
% Given a possibly high-level, ambiguous, or under-specified natural language instruction, a code agent should demonstrate capacity to \textit{discover context} and \textit{plan}, such as breaking down a high level objective into separate tasks that can be completed in a modular fashion. (navigating the task setting to gain a better understanding of the environment and identify relevant context.)
% Beyond correctness and efficiency, it may also be important that a model's task execution approaches align with human preferences (i.e. \textit{how} to do a task, not just \textit{what}) and can adjust its approach in response to evolving feedback.
% Across a longer task horizon that involves repetitive engagement with a common user or user group, an agent's ability to selectively apply historical dialogue and context can improve the efficiency and accuracy with which it carries out future tasks.
% In addition to exploring the capabilities of a single actor, \benchmark{} can also serve as a playground for evaluating multi-agent collaboration~\cite{havrylov2017emergence, wong2022deep}; such directions can draw heavily on existing software engineering paradigms for cooperation~\cite{dingsoyr2012decade, williams2000strengthening} to explore how teams of agents might be able to complete more significant engineering undertakings such as building and deploying a full stack application. \sy{make it more readable; can also cut if mentioned in sec 4}
% \section{The \benchmark{} Benchmark}
\label{sec:benchmark}

We present the \benchmark{} benchmark, a framework for designing interactive code environments. In this section, we first discuss the task formulation put forth by \benchmark{}. We then dive into the implementation of the \benchmark{} interface and how it enables the creation of efficient, reproducible, and configurable interactive coding tasks from existing and novel datasets. Finally, we demonstrate \benchmark{}'s potential to bring forth the variety of language grounding and understanding challenges that exist in the space of interactive coding.

\subsection{Formulation}
\label{sec:benchmark:formulation}
The \benchmark{} benchmark formalizes the interactive coding task as a partially observable Markov decision process (POMDP) with instruction space $\mathcal{U}$, state space $\mathcal{S}$, action space $\mathcal{A}$, state observation space $\mathcal{O}$, transition function $\mathcal{T}: \mathcal{S} \times \mathcal{A} \to \mathcal{S}$, and reward function $\mathcal{R}: \mathcal{S} \times \mathcal{A} \to \{0, 1\}$. Mirroring real world coding paradigms, the primary interactive setup presented by \benchmark{} is, given a natural language instruction $u \in \mathcal{U}$, an agent issues code as an action $a_t \in \mathcal{A}$ and receives execution output as a text observation $o_{t+1} \in \mathcal{O}$ across multiple turns. Depending on the nature of the task and system, the action space consists of 1) one or more programming languages and 2) the \texttt{submit} keyword. Given the constraints imposed by the environment, task episode, and execution context at a particular timestep, an action is deemed \textit{admissible} if can be parsed and executed within the environment ($\mathcal{T}(s, a) \neq s$)~\cite{yao2020calm}. Upon submission of the \texttt{submit} action at the state at timestep $T$, the task episode ends and a reward $r \in \mathcal{R}(s_T, \texttt{submit}) = \{0, 1\}$ is computed, where $1$ is awarded for successful completion of the task in its entirety, and $0$ for otherwise. We use the \textbf{Success Rate (SR)} metric, defined as the proportion of task episodes where $r = 1$. We also define the \textbf{Error \%} metric, which corresponds to the percentage of actions across a single task episode that is \textit{not} admissible.

A state $s_t \in \mathcal{S}$ generally represents the system or setting that the interactive coding task is taking place in (i.e. MySQL Database, Bourne Shell). Depending on the system, the state includes information such files/tables/folders, permissions, current working directory, execution context, and details about the OS or installed dependencies, to name a few. This latent state is not directly available to an agent, but can be uncovered via observations from agent actions. At a high level, the task formulation and feedback loop presented by the \benchmark{} task formulation reflects the development practices of programmers and software engineers.



\subsection{Design Choices}
\label{sec:benchmark:principles} \benchmark{} aims to serve as an expressive, safe, and lightweight framework for interactive code settings that subscribe to the aforementioned task formulation and can be used as a tool or training platform.
The \texttt{\benchmark{}Env} is the main deliverable towards this objective, providing foundational scaffolding for defining an interactive task environment that handles interaction efficiently and safely while also exposing a curated set of functions for task customization at numerous levels of granularity.

\textbf{Virtual Containers.} For developing lightweight and transferable interactive code settings with rich and diverse observation $\mathcal{O}$ and state $\mathcal{S}$ spaces, \benchmark{} leverages virtual containers, specifically Docker~\cite{merkel2014docker}, as a general execution sandbox.
Docker images, which can be custom made with a \texttt{Dockerfile}, provide excellent coverage of operating systems, languages, dependencies, and packages.
This affords task designers an exciting level of expressivity to manifest a wide range of task settings that can be accessed via a uniform, task-agnostic interface.
Execution in a virtual setting also provides safety against actions $\mathcal{A}$ that might be catastrophic on real world systems.
% \jy{Intercode makes it easy to set guardrails on the action space, so it serves as a good interface for interacting with real world systems too.}

% \textbf{Two Usage Modes.} \benchmark{} can run in two modes: \texttt{tool} and \texttt{eval}. In tool mode, \benchmark{}’s interface for interactive code settings can be used as a tool that language models can use to solve downstream tasks, or as a wrapper for connecting code agents to real world code tasks and settings. The inclusion of a dataset, which sets InterCode in \texttt{eval} mode, enables InterCode’s use as a training and testing environment for code and language agents. \sy{do we even use these terms? let's just say we use execution for both observation and reward}

\textbf{Customizable Execution, Dataset, Reward Function.} Within the interactive coding task formulation, \benchmark{} exposes key functionality that allows users to tailor an environment according a task's specific requirements and preferences. These include how an action is executed in a container (transition function $\mathcal{T}$), the criteria for task completion (reward function $\mathcal{R}$), and the dataset for training and evaluation (instruction space $\mathcal{U}$). A description of the \texttt{\benchmark{}Env} base class and methods for customization are fully described in \S\ref{appx:env:interface}.

\subsection{Usage of \benchmark{}}
\label{sec:benchmark:usage}
\sy{shorten? and tie to formulation terms}
\benchmark{} unifies the logic for facilitating interactive coding, and abstracts this functionality away such that what remains to implement are only task specific details.
To this end, creating an interactive coding task and environment using \benchmark{} requires task designers to provide three components.
\begin{itemize}[noitemsep,topsep=0pt]
    \item A \textit{Dockerfile} to define a containerized application for code execution.
    \item A \textit{dataset} containing pairs of (natural language query, reference gold command)
    \item A \textit{reward function} that evaluates an agent's actions against the reference gold command to determine task completion.
\end{itemize}

\textbf{Define Task Setting.} Given the name of a docker image, \benchmark{} handles creating and connecting to a corresponding virtual container that serves as a sand-boxed execution environment.
The container entity is also exposed as a variable to task designers in the case that custom logic is required to implement how to 1) Execute an action $a \in \mathcal{A}$  within the container and/or 2) Retrieve information from the container for reward calculation purposes.

Specifying a docker image from Docker Hub or defining a custom one with \texttt{Dockerfile} is the task designer's responsibility and creative outlet to select or design a suitable task setting, which can further be enriched with additional features (i.e. custom file system, more database tables, extra packages/dependencies, ) as needed for the task.
In this manner, a task setting can be configured without any impact on the downstream interaction logic.

\textbf{Define Dataset.} \benchmark{} requires that a dataset has at minimum two fields: \texttt{query}, a natural language instruction $u \in \mathcal{U}$, and \texttt{gold}, a code block that is a procedure for generating the correct answer. 
Additional fields can be included and incorporated for task-level customization via the \texttt{\benchmark{}Env} API.
Per task episode, the instruction text from an entry is shown directly or wrapped in a more elaborate prompt. Upon performing the \texttt{submit} action, the latest state and observation are graded against the \texttt{gold} command output.

This basic specification makes it easy to both conform existing datasets and devise new ones that fit \benchmark{}'s task formulation while offering flexibility for fabricating a wide range of coding task instructions.
In addition, the choice of <\texttt{query}, \texttt{gold}> pairs (as opposed to <\texttt{query}, \texttt{answer}>) encourages the construction of datasets that can be readily employed across different \benchmark{} task settings that may share the same programming language(s) for interaction such that the answer may vary across these settings. This presents opportunities to assess the generalizability, robustness, and adaptability of generations on top of completion correctness.

\textbf{Define Task Completion.} Across a single task episode, each interaction iteration of an action, observation, and state modification are implicitly logged by \benchmark{}.
At the \texttt{submit} time step, the simplest way to decide whether these sequence of actions complete the original task is with exact match comparison against the corresponding trajectory and outputs from executing the \texttt{gold} command.
With that said, given that the definition of completion can be heavily influenced by the nature of the instruction, task setting, and modes of interaction, exact match may often be too stringent of a rating style to use. 
Consequently, we make it possible to override exact match with custom reward functions using \texttt{\benchmark{}Env}, discussed further in \S\ref{appx:env:interface}.

\subsection{\benchmark{} Based Environments}
Following the procedure discussed in Section~\ref{sec:benchmark:usage}, we create two separate \benchmark{} based environments where Bash and SQL are the action spaces respectively. Table~\ref{table:env-compare} summarizes each environment.

\begin{table}
\centering
\begin{tabularx}{\textwidth}{lllll}
\toprule
    Action Space
    & Environment \inlinegraphics{figures/assets/emojis/docker.pdf}
    & Dataset \inlinegraphics{figures/assets/emojis/dataset.pdf}
    & Reward Function \inlinegraphics{figures/assets/emojis/reward.pdf} \\
\midrule
    Bash & Ubuntu Terminal & NL2Bash~\cite{lin-etal-2018-nl2bash} (200) & Latest Std. Output + File System $\Delta$ \\
    SQL & MySQL Database & Spider 1.0~\cite{yu-etal-2018-spider} (1034) & Latest Std. Output \\
    Python & Python Interpreter & MBPP~\cite{austin2021program} (117)   & Submitted Function \\
\bottomrule
\end{tabularx}
\vspace{0.5em}
\caption{Rundown of the two environments with Bash and SQL as action spaces developed using the \benchmark{} framework. The numbers in parentheses refer to the number of task instances adopted from each dataset. Each environment is defined in under 200 lines of code total. Specific discussion of the environment construction and reward function can be found in \S~\ref{appx:env:bashenv} and \S~\ref{appx:env:sqlenv}}
\label{table:env-compare}
\vspace{-1.5em}
\end{table}

\textbf{Bash Environment.} Figure~\ref{fig:preview} visualize the interactive flow and evaluation procedure for the \benchmark{}-Bash environment. We define a \texttt{bash} shell within an Ubuntu Operating System as the task setting.
To evaluate an agent's ability to adapt generations to different situations, we architect four distinct file systems that can be swapped into the Bash environment by changing a single line in the \texttt{Dockerfile}.
We bootstrap the NL2Bash~\cite{lin-etal-2018-nl2bash} dataset to create an interactive coding task where an agent completes an instruction via \texttt{bash} actions. Transferring NL2Bash to the interactive task setting requires simple transformations to ground instructions and \texttt{gold} code blocks in the file system.
Finally, we define a custom reward function that evaluates an agent's performance against file system modifications and the latest execution output.
For instance, an instruction asking an agent to write content to a text file then get the file's new size in human readable format would check that both 1. The text file was modified correctly (i.e. \texttt{echo "Hello" $>>$ file.txt} and 2. The file's updated size is shown in the latest execution output (i.e. \texttt{ls -lh file.txt})
More specific details are in \S~\ref{appx:env:bashenv}.

\textbf{SQL Environment.} We define an SQL interpreter within MySQL database as the task setting, then similarly adapt the Spider~\cite{yu-etal-2018-spider} dataset to create an interactive coding task where an agent answers questions in the form of instructions with a set of records as execution output from \texttt{SQL} actions. We perform an intersection over union comparison of the latest execution outputs of the agent's commands and gold command(s) to determine task completion. Visualizations like Figure~\ref{fig:preview} for SQL along with a more extensive implementation discussion for this environment are in \S~\ref{appx:env:sqlenv}

\subsection{Research Challenges Imposed by \benchmark{}}
\label{sec:benchmark:challenges}
\benchmark{}'s interface and design provides fertile ground for constructing language understanding and decision making challenges spanning across NLP and RL.
Solving a task across multiple iterations implicitly relies on a model's ability to perform \textit{error correction} by determining what action to perform next based on past execution output(s).
Given a possibly high-level, ambiguous, or under-specified natural language instruction, a code agent should demonstrate capacity for \textit{planning}, such as navigating the task setting to gain a better understanding of the environment and identify relevant context.
Such \textit{reasoning} is also useful for breaking down a high level objective into separate tasks that can then be completed in a modular fashion.
Beyond correctness and efficiency, it may also be important that a model's task execution approaches align with human preferences (i.e. \textit{how} to do a task, not just \textit{what}) and can adjust its approach in response to evolving feedback. Across a longer task horizon that involves repetitive engagement with a common user or user group, an agent's ability to selectively apply historical dialogue and context can improve the efficiency and accuracy with which it carries out future tasks. In addition to exploring the capabilities of a single actor, \benchmark{} can also serve as a playground for evaluating multi-agent collaboration~\cite{havrylov2017emergence, wong2022deep}; such directions can draw heavily on existing software engineering paradigms for cooperation~\cite{dingsoyr2012decade, williams2000strengthening} to explore how teams of agents might be able to complete more significant engineering undertakings such as building and deploying a full stack application. \sy{make it more readable; can also cut if mentioned in sec 4}
\section{Methods}
\label{sec:methods}
We perform preliminary experiments to gauge the proficiency and behavior of current large language models on interactive coding tasks with Bash and SQL.
To observe and elicit relevant reasoning skills, we draw on several existing prompting strategies that have been put forth to augment language models' reasoning and problem-solving skills.
We apply these prompting strategies to models across the following three families: OpenAI (\texttt{text-davinci-003}, \texttt{gpt-3.5-turbo}, \texttt{gpt-4}), PaLM-2 (\texttt{text-bison-001}, \texttt{chat-bison-001}) \cite{anil2023palm}, and Open Source (Vicuna-13B \cite{vicuna2023}, StarChat-16B \cite{li2023starcoder}).

% \subsection{Prompting strategies}
Figure~\ref{fig:prompting-approaches} visualizes the four adjusted prompting strategies we evaluate on \benchmark{}.

\textbf{Single Turn} is a zero-shot attempt. A model is given a simple description of the task setting and asked to generate code in a specific programming language that would address the query. The first generation in response to the user's question is then evaluated in the \benchmark{} environment.

\begin{figure}
    \centering
    \includegraphics[width=1.0\textwidth]{figures/assets/Approaches.pdf}
    \caption{Overview of Prompting Strategies adjusted for evaluation on \benchmark{}. The "Try Again" termination constraint is conditioned on reward = 1, while ReAct~\cite{yao2023react} and Plan \& Solve~\cite{wang2023planandsolve} are determined by the agent itself. This is because the purpose of the "Try Again" method is to explore how capable agents are at error correction from feedback, while the other two are more concerned with the overall success of general problem-solving strategies.}
    \label{fig:prompting-approaches}
\end{figure}

\textbf{"Try Again"} is an iterative feedback set up. In the initial message, the agent is informed of the task setting and its interactive nature; an agent has multiple turns to interact with the system, wherein each turn, upon generating an action, the execution output of the action is fed back as an observation. This continues until a reward of 1 (task completion) is achieved or the number of turns ($n$) is exhausted. The agent's position in this approach is meant to mirror human software development as closely as possible. The goal of this method is to probe language models' raw interactive coding abilities in addition to illustrating the benefits and different challenges that arise in interactive coding tasks.

\textbf{ReAct and Plan \& Solve.} We write prompts and design workflows that follow the text and task configurations described in ReAct~\cite{yao2023react} and Plan \& Solve~\cite{wang2023planandsolve} as faithfully as possible. For these two approaches, the termination of a task episode is conditioned upon the agent's own judgment, as our goal with these methods is to gauge the transferability to and efficacy of existing reasoning frameworks with respect to the interactive coding task.
Full prompt templates are included in \S\ref{appx:experiments:templates}.

\section{Experiments}
\label{sec:experiments}
% In our experiments, we first examine both the benefits of an interactive setting towards coding tasks along with the gaps and research challenges it presents that models have yet to overcome.
% We then demonstrate how the interactive coding task typically encapsulates a multitude of reasoning challenges, motivating a higher order question of how models might go about exercising and alternating across multiple problem solving strategies on its own volition.
% Finally, we present a small-scale dataset inspired by Capture the Flag, a popular genre of computer security competitions, to showcase \benchmark{}'s value as a testbed for new datasets and task formats in the realm of interactive code.
% Specific trajectories that are evidence of the conclusions developed across the following sections can be found in \S~\ref{appx:experiments:trajectory_analyses}.

\begin{table}
\centering
\begin{tabular}{llllll|lllll}
\toprule
      \benchmark{}-SQL & \multicolumn{5}{c}{Single Turn} & \multicolumn{5}{c}{Try Again ($n$ = 10)} \\
     Model / Hardness & Easy & Med & Hard & Extra & All & Easy & Med & Hard & Extra & All \\
\midrule
text-davinci-003 & 20.6 & 4.9 & 1.7 & 0.0 & 7.4 & 32.4 & 14.6 & 5.2 & 4.2 & 15.6 \\
gpt-3.5-turbo & 22.6 & 8.3 & \textbf{5.7} & \textbf{3.6} & 10.5 & 72.5 & 44.3 & 43.7 & 21.1 & 47.3 \\
gpt-4 & 19.8 & 7.2 & 4.6 & 3.0 & 9.1 & \textbf{87.5} & \textbf{76.7} & \textbf{66.7} & \textbf{52.4} & \textbf{73.7} \\
text-bison-001 & \textbf{23.8} & \textbf{10.9} & \textbf{5.7} & 0.6 & \textbf{11.5} & 27.0 & 12.3 & 5.7 & 0.6 & 12.9 \\
chat-bison-001 & 18.5 & 6.5 & 4.0 & 0.0 & 7.9 & 22.2 & 7.8 & 6.9 & 0.0 & 9.9 \\
% Claude-v1 & & & & & & & & & & \\
Vicuna-13B & 8.1 & 1.3 & 0.6 & 0.0 & 2.6 & 18.9 & 3.4 & 1.7 & 0.0 & 6.3 \\
StarChat-16B & 21.8 & 7.4 & 2.9 & 0.0 & 8.9 & 22.3 & 8.5 & 2.9 & 1.2 & 9.7 \\
\bottomrule
\end{tabular}
\vspace{5pt}
\caption{Success Rate for single vs.\,multi turn evaluation on \benchmark{}-SQL (refer \S\ref{appx:env:sqlenv}). Query difficulty is adopted from Spider~\cite{yu-etal-2018-spider}. Best metrics are in \textbf{bold}.}
\label{table:single-vs-multi-sql}
\vspace{-10pt}
\end{table}
\begin{table}
\centering
\begin{tabular}{llllll|lllll}
\toprule
     \benchmark{}-Bash & \multicolumn{5}{c}{Single Turn} & \multicolumn{5}{c}{Try Again ($n$ = 10)} \\
Model / File System & 1 & 2 & 3 & 4 & All & 1 & 2 & 3 & 4 & All \\
\midrule
text-davinci-003 & 10.0 & 32.1 & 28.8 & 33.3 & 24.6 & 30.0 & \textbf{52.8} & 32.2 & 44.4 & 38.7 \\
gpt-3.5-turbo & \textbf{30.0} & \textbf{39.6} & 33.3 & 37.0 & \textbf{34.5} & \textbf{45.0} & 49.1 & 45.0 & 48.1 & 46.5 \\
gpt-4 & 25.0 & 37.7 & \textbf{36.7} & \textbf{40.7} & 34.0 & 41.7 & 47.2 & \textbf{51.7} & \textbf{59.2} & \textbf{48.5} \\
text-bison-001 & 15.0 & 22.6 & 11.7 & 22.2 & 17.0 & 23.3 & 28.3 & 16.7 & 22.2 & 22.5 \\
chat-bison-001 & 12.1 & 22.5 & 16.7 & 22.2 & 17.7 & 13.8 & 24.5 & 18.3 & 22.2 & 19.2 \\
% Claude-v1 & & & & & & & & \\
Vicuna-13B & 10.0 & 24.5 & 18.3 & 7.4 & 16.0 & 15.0 & 35.8 & 25.0  & 22.2 & 24.5 \\
StarChat-16B & 15.5 & 22.6 & 13.3 & 22.2 & 17.7 & 17.2 & 30.2 & 21.7 & 29.6 & 23.7 \\
\bottomrule
\end{tabular}
\vspace{5pt}
\caption{Success Rate across file systems for single vs. multi-turn evaluation on  \benchmark{}-Bash (refer \S\ref{appx:env:bashenv}). To evaluate models' ability to interact with different task settings, we evaluate disjoint sets of Bash instructions across four different file systems. Best metrics are in \textbf{bold}.}
\label{table:single-vs-multi-bash}
\vspace{-1.5em}
\end{table}

\subsection{Base models comparison}
\label{sec:experiments:base_compare}
\textbf{Task performances.} We first compare the success rate of models in the Single Turn and Try Again settings for both the \benchmark{}-Bash and SQL datasets.
From Table~\ref{table:single-vs-multi-sql} and Table~\ref{table:single-vs-multi-bash}, we observe that performance across different levels of task difficulty (SQL) and different file systems (Bash) is superior in the interactive setting for all models, with a notable multi-fold increase for GPT-4 ($9.1\% \rightarrow 73.7\%$) on the \benchmark{}-SQL task.

\textbf{Analysis of interactions.} Manual inspection of trajectory logs indicates that models actively exercise later turns for discovering relevant context, correcting errors via execution feedback as observations, and solving problems via iteratively constructing and editing actions as affirmed by Figure~\ref{fig:turn_stats}.
In addition, models also demonstrate a level of planning and modular problem solving; for instructions with \texttt{gold} commands that chain multiple commands together (i.e. with \text{\textbar}, $>$, or \texttt{;} in \texttt{bash}) or consist of multiple sub-problems (i.e. subqueries in \texttt{SQL}), models will use observations from solving smaller sub-problems in earlier turns to compose the higher-order action. Trajectories that exhibit these phenomena are in \S~\ref{appx:experiments:trajectory_analyses}

\textbf{Failure cases.} With that said, both Figure~\ref{fig:turn_stats} exhibits a plateauing in Success Rate and and Error \%.
This suggests that as the amount of context and feedback builds up, models are less capable of discerning relevant past history toward future actions.
In late-turn scenarios, task episode trajectories often reveal repetition of earlier actions, a failure to effectively use recent observations towards deciding an appropriate next action, or an inability to recognize that a current problem-solving chain of thought is inconclusive or futile.
This is particularly evident for \texttt{hard} and \texttt{extra} level \benchmark{}-SQL task instructions that require context spanning across several tables and actions that incorporate multiple clauses.
We note that even when the full schema of all tables and their descriptions are offered in addition to the original instructions, models still benefit greatly from using interaction to experiment with different \texttt{JOIN} and filtering operators across multiple turns, as demonstrated in \S~\ref{appx:experiments:additional_expr}.
A larger context window size, retrieval of useful memory, and more adaptive reasoning paradigms are just a handful of potential solutions to overcoming such challenges.

\subsection{Prompting strategy comparison}
\label{sec:experiments:strategy_compare}
Initiating language agents with prompting strategies that encourage different forms of reasoning toward problem-solving improves performance on the interactive coding task to varying degrees. Table~\ref{table:strategy-comparison} presents side-by-side comparisons of the success rate, number of turns, and error rate per strategy. 
Compared to Try Again, which lacks specific guidance on leveraging multiple turns, more explicit reasoning frameworks such as ReAct and Plan \& Solve policies generally achieve higher success rates (SQL: $47.3\% \rightarrow 58.7\%$) with fewer turns and a higher rate of admissible commands.

\begin{figure}
    \begin{subfigure}[t]{0.5\textwidth}
        \centering
        \includegraphics[width=\textwidth]{figures/assets/bash_turn_stats.pdf}
        \caption{Success rate vs. turns for \benchmark{}-Bash}
    \end{subfigure}
     \begin{subfigure}[t]{0.5\textwidth}
     \centering
        \includegraphics[width=\textwidth]{figures/assets/sql_turn_stats.pdf}
        \caption{Success rate vs. turns for \benchmark{}-SQL}
    \end{subfigure}
    \caption{Growth in Success Rate with increase in number of interaction turns across models configured with Try Again prompting strategy for \benchmark{}-Bash and SQL tasks.}
    \label{fig:turn_stats}
\end{figure}
\begin{table}
\centering
\begin{tabular}{llll|lll|lll}
\toprule
     & \multicolumn{3}{c}{Try Again ($n$ = 10)} & \multicolumn{3}{c}{ReAct ($n$ = 10)} & \multicolumn{3}{c}{Plan \& Solve} \\
     & SR & Turns & Error \%
     & SR & Turns & Error \%
     & SR & Turns & Error \% \\
\midrule
SQL & 47.3 & 7.25 & 46.4
    & \textbf{58.7} & 5.30 & \textbf{6.94}
    & 49.1 & \textbf{4.29} & 16.2 \\
Bash & \textbf{46.5} & 6.15 & 24.9
    & 20.5 & \textbf{4.40} & \textbf{20.4}
    & 28.0 & 6.65 & 53.3 \\
\bottomrule
\end{tabular}
\vspace{5pt}
\caption{Comparison of different prompting strategies across the entire \benchmark{}-SQL and \benchmark{}-Bash datasets using \texttt{gpt-3.5-turbo} as the base model. \textit{Turns} refers to the average number of turns taken for a single task episode. For Try Again and ReAct, the max number of turns $n = 10$. The highest Success Rate, fewest Turns, and lowest Error \% are highlighted per dataset since they reflect more accuracy and efficient task solving. Best metrics are in \textbf{bold}.}
\label{table:strategy-comparison}
\vspace{-1.5em}
\end{table}

\textbf{Different tasks present different learning challenges.}
An important skill to solving the \benchmark{}-SQL task is the ability to discover context and construct actions conditionally based on information revealed in prior observations.
Given that \benchmark{}-SQL task instructions are phrased most commonly as questions, adapting to the task setting and new information discovered along the way puts more emphasis on error correction and context discovery.
On the other hand, the more declarative and multi-step nature of the \benchmark{}-Bash task instructions is more aptly solved by planning and modular task completion.
These distinctions manifest in the Plan \& Solve strategy's performance gap between the \benchmark{}-SQL and \benchmark{}-Bash tasks; while Plan \& Solve encourages a model to decompose problems into more manageable steps, the strategy is less favorable towards adjusting on the fly in response to execution feedback.
Example trajectories supporting these claims are in \S~\ref{appx:experiments:trajectory_analyses}.

\textbf{More adaptive reasoning is favorable.}
Compared to "imperative" reasoning paradigms such as Plan \& Solve which prescribe a relatively rigid procedure, more flexible frameworks like ReAct, which do not enforce any particular logical formula or roadmap, are more conducive to eliciting a broader set of reasoning capabilities.
However, while ReAct's performance is generally superior to Plan \& Solve, tasks solved by \textit{both} strategies with \texttt{gpt-3.5-turbo} make up $57\%$ ($407/708$) and $27.6\%$ ($21/76$) of the union of all successfully solved \benchmark{}-SQL and \benchmark{}-Bash tasks respectively.
This discrepancy highlights a trade-off between the guidance and structural constraints that are inherent to prompting strategies; schemes that draw out specific reasoning patterns often overlook other equally useful capabilities.
\benchmark{}'s interactive coding task can serve as a strong litmus test toward more adaptable, variegated model reasoning.

\subsection{New tasks \& datasets opportunities}
\benchmark{}'s task formulation, modular design, flexible task construction, and use of virtual containers enable task designers to manifest new, complex, code-driven tasks, where completion is much more attainable through interaction.
We draw inspiration from Capture the Flag (CTF)~\cite{cowan2003ctf}, a competitive cybersecurity game that requires expertise in coding, cryptography (i.e. binary exploitation, forensics), reverse engineering, and recognizing security vulnerabilities to accomplish the primary objective of discovering encrypted "flags" concealed within code snippets or file systems.
Compared to \benchmark{}-Bash \& -SQL, CTF is much more complicated, requiring an agent to exercise knowledge of multiple coding languages, modularize a higher-order objective into sub-problems, construct multi-step plans towards solving each problem, and adjust strategy when a plan fails to yield any useful insights.

We establish \benchmark{}-CTF, a new dataset consisting of 100 CTF objectives from picoCTF~\cite{picoctf}.
Following the interactive coding task formulation, each task instance in \benchmark{}-CTF is given as a \texttt{<instruction, assets, hidden flag>} tuple.
We first construct a Bourne Shell within an Ubuntu OS as the task environment.
Here, \benchmark{}'s use of virtual containers is crucial, as necessary actions can be irreversibly damaging on real systems (i.e. \texttt{rm -rf}, \texttt{sudo} access).
Per task instance, the associated assets (e.g., images, executables, code), necessary for task completion, are copied into the OS file system.
Given this setting, a task worker must understand the given material and investigate the assets to develop potential solutions.
Executing a successful approach must be done across multiple steps with various conditionals, where the execution feedback of a prior step could have a significant effect on the next step.
% For more complex problems, some solution result in an inconclusive dead end.
% In these situations, models must not only recognize when they should pivot their approach, but also decide what approach makes the most sense to attempt given prior failures.
Figure~\ref{fig:traj_ctf} spotlights the diverse skills needed for CTF. 

\begin{figure}
    \centering
    \makebox[\textwidth][c]{\includegraphics[width=\textwidth]{figures/assets/Traj-ctf-forensics-wide.pdf}}
    \caption{GPT-4's interaction trajectory for a binary exploitation CTF task. This requires proficiency in \textcolor{blue}{Bash} and \textcolor{OliveGreen}{Python}, among additional knowledge and reasoning. \textcolor{orange}{Orange} text and arrows highlight the feedback that the model attends to in generating the next action. In last step, agent submits flag.}
    \label{fig:traj_ctf}
\end{figure}

%%%% MARK: NON INCLUDED SECTIONS %%%%

% \subsection{Intercode Tasks Present Opportunities to design better techniques for soliciting different/multiple capabilities from LLMs}
% \jy{This message might not be worth explaining or the amount of explanation it would require (1. explain plan\&solve + refine and why it was made, 2. Discuss how it performs better than other strategies 3. Still not perfect)}
% \input{figures/tables/expr-plan-solve-refine}

% \subsection{Tools / Sim-to-Real Use Cases}
% Nothing actually done in this area, but is this a worthwhile thing to have + can we do something in the next two weeks? Can be a toy thing.

% IDEA: Take existing code and refactor it to be (more readable/efficient, handle some errant edge case). We don't have unit tests - use an environment to interactively verify that the refactored version works the same as before + achieves the improvement. (Requires Web + Bash/SQL?)
\section{Discussion}
\label{sec:conclusion}

\textbf{Conclusion.} We have developed \benchmark{}, a novel lightweight framework that facilitates interaction between Language Models and the underlying environment, enabling them to mimic the human approach to language-to-code generation. Our framework has shown promising results when applied to state-of-the-art models using different prompting styles. It effectively leverages the capabilities of LMs to break down complex tasks and recover from errors within a secure and isolated environment. The ability to seamlessly convert existing datasets into the interactive format using \texttt{\benchmark{}Env} API, and furthermore, the Bash and SQL environments, empowers task designers to construct new tasks to unlock the plethora of challenges that await in the space of interactive coding.

\textbf{Limitations and future directions.} We point out several current limitations of \benchmark{}. At this time, the number of \benchmark{} based environments is limited to Bash, SQL, and Python action spaces and datasets; within the near future, we plan to expand the number of offerings to cover a wider set of programming languages and datasets that should further deliver on \benchmark{}'s purported promises of efficient and expressive task construction. Second, the CTF dataset is limited to just four task instances due to our manual curation procedure. We hope to release more formal work soon that provides a more thorough analysis of the reasoning and collaboration challenges of the CTF task along with a more extensive dataset for evaluation purposes.
\newpage
\section*{Acknowledgements}
We thank Xiao Liu for the Vicuna/Alpaca APIs, Carlos Jimenez and Yuhan Liu for trying our code, and Princeton NLP Group for helpful discussion and feedback in general. We acknowledge support from the National Science Foundation under Grant No. 2107048. Any opinions, findings, and conclusions or recommendations expressed in this material are those of the author(s) and do not necessarily reflect the views of the National Science Foundation.
\label{sec:acknowledgements}
% \newpage
\bibliography{main}
\bibliographystyle{abbrvnat}
% \section*{Checklist}

\begin{enumerate}

\item For all authors...
\begin{enumerate}
  \item Do the main claims made in the abstract and introduction accurately reflect the paper's contributions and scope?
    \answerYes{}
  \item Did you describe the limitations of your work?
    \answerYes{} (See \S\ref{appx:misc:potential_impacts})
  \item Did you discuss any potential negative societal impacts of your work?
    \answerYes{} (See \S\ref{appx:misc:potential_impacts})
  \item Have you read the ethics review guidelines and ensured that your paper conforms to them?
    \answerYes{}
\end{enumerate}

\item If you ran experiments (e.g. for benchmarks)...
\begin{enumerate}
  \item Did you include the code, data, and instructions needed to reproduce the main experimental results (either in the supplemental material or as a URL)?
    \answerYes{} (See beginning of the appendix)
  \item Did you specify all the training details (e.g., data splits, hyperparameters, how they were chosen)?
    \answerYes{} (See \S\ref{appx:experiments:details}
  \item Did you report error bars (e.g., with respect to the random seed after running experiments multiple times)?
    \answerYes{} (See \S\ref{appx:experiments:additional_expr})
  \item Did you include the total amount of compute and the type of resources used (e.g., type of GPUs, internal cluster, or cloud provider)?
    \answerYes{} (See \S\ref{appx:experiments:details})
\end{enumerate}

\item If you are using existing assets (e.g., code, data, models) or curating/releasing new assets...
\begin{enumerate}
  \item If your work uses existing assets, did you cite the creators?
    \answerYes{}
  \item Did you mention the license of the assets?
    \answerYes{} (See \S\ref{appx:env:bashenv} and \S\ref{appx:env:sqlenv})
  \item Did you include any new assets either in the supplemental material or as a URL?
    \answerYes{} (See beginning of the appendix)
  \item Did you discuss whether and how consent was obtained from people whose data you're using/curating?
    \answerYes{} (See \S\ref{appx:experiments:human_perf})
  \item Did you discuss whether the data you are using/curating contains personally identifiable information or offensive content?
    \answerYes{} (See \S\ref{appx:misc:data_collection_cklist})
\end{enumerate}

\item If you used crowdsourcing or conducted research with human subjects...
\begin{enumerate}
  \item Did you include the full text of instructions given to participants and screenshots, if applicable?
    \answerYes{} (See \S\ref{appx:experiments:human_perf})
  \item Did you describe any potential participant risks, with links to Institutional Review Board (IRB) approvals, if applicable?
    \answerYes{} (See \S\ref{appx:experiments:human_perf})
  \item Did you include the estimated hourly wage paid to participants and the total amount spent on participant compensation?
    \answerYes{} (See \S\ref{appx:experiments:human_perf})
\end{enumerate}

\end{enumerate}
\newpage
\appendix
\section*{Appendix}
\label{appx:intro}
In this appendix, we provide additional details about the implementation and usage of the \benchmark{} framework and the \texttt{\benchmark{}Env} interface. We also provide visualizations and analyses of additional experiments to demonstrate \benchmark{}'s utility and garner further insight into the extent of current models' performance on the interactive coding task. The full template for each prompting strategy is also included. Finally, we also discuss some of the impacts, risks, and limitations of our work. The webpage for \benchmark{} is \url{https://intercode-benchmark.github.io/}. The code for \benchmark{} is \url{https://github.com/princeton-nlp/intercode}; the link is also included on the \benchmark{} webpage.

\section{Environment Details}
\label{appx:env}

\subsection{\benchmark{} Interface}
\label{appx:env:interface}
The \benchmark{} interface inherits the OpenAI gym~\cite{openai_gym} environment API definition. 
Specifically, \texttt{\benchmark{}Env} is written as an abstract class that primarily handles the main execution logic for processing code interactions, in addition to logging, data management, and sand-boxed execution, along with both environment-level and task-level customization. 

\texttt{\benchmark{}Env} exposes the following API. Creating an interactive coding environment requires defining a subclass of \texttt{\benchmark{}Env}. The methods denoted with an asterisk can be overridden for the purposes of customization.

\texttt{\_\_init\_\_(self, data\_path: str, image\_name: str, **kwargs)}
\begin{itemize}[noitemsep,topsep=0pt]
    \item Validates that the dataset specified by \texttt{data\_path} is formatted correctly and can be used in an interactive setting.
    \item Uses the Docker image specified by \texttt{image\_name} to create and connect with a Docker container instance of the image.
    \item Initializes Logging Handler
    \item Keyword arguments:
    \begin{itemize}[noitemsep]
        \item \texttt{verbose (bool)}: If true, logging is enabled and environment interactions are shown to standard output
        \item \texttt{traj\_dir (str)}: If a valid path is provided, task episode summaries are saved to the given directory (generated by \texttt{save\_trajectory})
        \item \texttt{preprocess (callable)}: If provided, this function is run before every task episode. It is a way to provide task instance-specific customization of the execution environment.
    \end{itemize}
\end{itemize}

\texttt{reset(self, index: int = None) -> Tuple[str, Dict]}
\begin{itemize}[noitemsep,topsep=0pt]
    \item Retrieves task record from data loader
    \item Calls \texttt{reset\_container}
    \item Reset task level logger, instance variables
\end{itemize}

\texttt{step(self, action: str) -> Tuple[str, int, bool, Dict]}
\begin{itemize}[noitemsep,topsep=0pt]
    \item Log (action, observation)
    \item Invoke \texttt{exec\_action} on action argument
    \item If \texttt{action}=\texttt{submit}, invoke \texttt{get\_reward}, \texttt{save\_trajectory}
\end{itemize}

\texttt{save\_trajectory(self)}
\begin{itemize}[noitemsep,topsep=0pt]
    \item Saves task metadata, (action, obs.) sequence, and reward info to \texttt{.json} in \texttt{traj\_dir}
\end{itemize}

\texttt{close(self)}
\begin{itemize}[noitemsep,topsep=0pt]
    \item Safely exit or stop any resources (i.e. docker container) used by the environment
\end{itemize}

\texttt{* execute\_action(self, action: str)}
\begin{itemize}[noitemsep,topsep=0pt]
    \item Defines how the action is executed within the context of the docker container.
    \item Requires impl. because the Dockerfile definition, particularly its \textit{entrypoint}, affects how an action would be invoked within the container.
    \item Default impl. passes the \texttt{action} string directly into a \texttt{self.container.exec(action)} call, which invokes the action in the environment and returns execution output. A timeout is imposed on execution duration.
\end{itemize}

\texttt{* get\_reward(self) -> Tuple[float, Dict]}
\begin{itemize}[noitemsep,topsep=0pt]
    \item Handles reward calculation of actions with respect to the gold command(s) for a task episode.
    \item Requires impl. because the concept and scoring for task completion varies across datasets and environments.
\end{itemize}

\texttt{* reset\_container(self)}
\begin{itemize}[noitemsep,topsep=0pt]
    \item Handles resetting of execution container (i.e. resetting file system to original state).
    \item Requires impl. because the approach to restoring a setting to its initial state varies.
\end{itemize}

\begin{figure}
    \centering
    \includegraphics[width=0.85\textwidth]{figures/assets/interface-visual.pdf}
    \caption{Visualization demonstrating the intended invocations and usage of the \texttt{\benchmark{}Env} interface, along with how the functions requiring implementation (\texttt{get\_reward()}, \texttt{execute\_action()}, \texttt{reset\_container()} are called by the methods of the main interactive loop.}
    \label{appx:fig:interface_visual}
\end{figure}

Figure~\ref{appx:fig:interface_visual} conveys how each of these methods are invoked and how they related to one another. In summary, the technicalities for setting up an interactive coding task for a specific system with one or more programming languages as the action space involve:

\begin{itemize}[noitemsep,topsep=0pt]
    \item Defining a \texttt{Dockerfile}
    \item Providing a dataset with the \texttt{query} and \texttt{gold} fields
    \item (Optional) Defining a reward (\texttt{get\_reward}) function to define task completion.
    \item (Optional) Creating an \texttt{\benchmark{}Env} subclass that overrides the \texttt{execute\_action} and \texttt{get\_reward} methods
\end{itemize}

\subsection{Bash Environment}
\label{appx:env:bashenv}
\textbf{Environment definition.} The \texttt{Dockerfile} defining the Bash-based environment is founded on the LTS version of the Ubuntu operating system. Several Linux dependencies that can potentially be used by an agent to address instructions in the \benchmark{}-Bash Dataset are then installed via the Advanced Package Tool (\texttt{apt}) interface. Next, a shell script is invoked within the \texttt{Dockerfile} to initialize one of the three file systems displayed in Figure~\ref{fig:nl2b_filesystems}. The shell script consists of a simple sequence of \texttt{mkdir}, \texttt{touch}, and \texttt{echo} commands to deterministically create and populate the content of multiple files and folders. Finally, \texttt{git} is configured for the purposes of determining file diffs per task episode (\texttt{git status -s}) and resetting an environment to its original state (\texttt{git reset --hard; git clean -fd;}) before the beginning of a new task episode. The original code for the \texttt{Dockerfile} along with the file system creation scripts can be found on the project GitHub repository.

% \textbf{\benchmark{}-Bash Dataset.} The \benchmark-Bash dataset is built over NL2Bash \cite{lin-etal-2018-nl2bash} which lacks specificity in queries and grounding to any underlying file system, hindering it from being directly used for interactive evaluations. 
% To overcome this limitation, we process NL2Bash as follows:
% \begin{itemize}[noitemsep,topsep=0pt]
% \item Consider a subset of 1000 commands with each having $\geq$ 4 utilities
% \item Filter out commands with unsupported utilities (eg. "ssh", "sudo", time and GUI dependent utilities)
% \item Filter out non-UNIX, non-Linux commands
% \item Enhance under-specified commands with specific file names/directory names/paths
% \item Update deprecated utilities/flags
% \end{itemize}
% The resulting 200 commands were grouped into 4 disjoint sets, 3 of which were grounded to custom-designed file systems, while one set is file-system agnostic. This categorization allowed for a comprehensive evaluation of different command-grounding scenarios.


\textbf{Dataset details.} The log-frequency distribution of the top-50 utilities is displayed in Figure~\ref{fig:nl2b_utility_distribution}.
The NL2Bash~\cite{lin-etal-2018-nl2bash} dataset is made available for use under the \hyperlink{https://www.gnu.org/licenses/gpl-3.0}{GPLv3 License}.
To assess the generalizability of our approach, we designed three distinct file systems to accommodate the bash commands we collected. A key consideration during the construction of these file systems was to ensure that a significant portion of the executed commands would not result in operations that yield no changes. This deliberate design choice aimed to provide a more comprehensive assessment of our approach's adaptability and effectiveness across various scenarios and command executions. The file systems encompass a wide range of file types, including text files (.txt), program files (.c, .java, .py), compressed files (.gz), shell scripts (.sh), PHP scripts (.php), JSON files (.json), documents (.doc), spreadsheets (.csv), webpages (.html), database schemas (.sql), hidden files, and files with special characters in their names, convoluted folder hierarchies. Their directory structures are illustrated in Figure~\ref{fig:nl2b_filesystems}. For simplicity, we consider the top-level folder created within the root directory (\texttt{testbed}, \texttt{system}, \texttt{workspace}) as the root of each file system. This root folder contains files and sub-folders that necessitate access and manipulation, while changes are monitored throughout the entire container to accurately evaluate the models' actions. Notably, we intentionally designed file system 1 to be more intricate and encompass relatively challenging bash tasks compared to the other two file systems. Thereby, the models' performance is relatively lower for file system 1.

\begin{figure}
\begin{subfigure}[t]{0.25\textwidth}
\dirtree{%
.1 /testbed.
.2 FooBar.html.
.2 Hello.java.
.2 Hello1.java.
.2 NewClass.java.
.2 dir1.
.3 AnotherHello.java.
.3 info.php.
.3 subdir1.
.4 jsonfile1.json.
.4 pythonscript4.py.
.4 shellscript1.sh.
.4 subsubdir1.
.5 pythonscript1.py.
.5 shellscript4.sh.
.4 textfile4.txt.
.3 subdir2.
.3 textfile1.txt.
.2 dir2.
.3 shellscript2.sh.
.3 subdir1.
.4 javafile1.java.
.4 textfile2.txt.
.3 subdir2.
.4 pythonscript2.py.
.4 shellscript5.sh.
.4 subsubdir1.
.4 textfile5.txt.
.2 dir3.
.3 subdir1.
.4 pythonscript3.py.
.4 subsubdir1.
.5 FooBar.
.6 file.txt.
.5 file.txt.
.5 shellscript3.sh.
.5 textfile3.txt.
.5 tmp.
.6 tmp.txt.
.3 subdir2.
.4 csvfile1.csv.
.3 textfile6.txt.
.2 hello.c.
.2 hello.php.
.2 recent.txt.
.2 textfile7.txt.
}
\caption{File System 1}
\label{fig:nl2b_fs_1}
\end{subfigure}
\hfill
\hfill
\begin{subfigure}[t]{0.25\textwidth}
\dirtree{%
.1 /system.
.2 .DS\_Store.
.2 MANIFEST.
.2 a.out.
.2 folder1.
.3 a.out.
.3 data.csv.
.3 doc1.doc.
.3 doc2.doc.
.3 keep.txt.
.3 log1.log.
.3 new.sh.
.3 old2.txt.
.3 recent.txt.
.3 script1.sh.
.3 text2.txt.
.3 text3.txt.
.3 text4.txt.
.2 folder2.
.3 special text3.txt.
.3 special\_text1.txt.
.3 special\_text2.txt.
.3 text1.txt.
.2 folder2.tar.gz.
.2 folder3.
.3 backup\_dbg.
.4 backup.
.4 sql1.sql.
.4 text1\_dbg.txt.
.3 special text4.txt.
.3 temp.
.4 empty.txt.
.4 temp1.
.5 temp.
.5 temp\_1.
.5 text1.txt.
.2 html1.html.
.2 temp.
.3 keep.txt.
.2 text1.txt.
.2 text3.txt.
}
\caption{File System 2}
\label{fig:nl2b_fs_2}
\end{subfigure}
\hfill
\begin{subfigure}[t]{0.25\textwidth}
\dirtree{%
.1 /workspace.
.2 .hidden.txt.
.2 archive.tar.gz.
.2 dir1.
.3 .hidden1.txt.
.3 a.txt.
.3 all.txt.
.3 file.c.
.3 file.txt.
.3 hello.c.
.3 hello.txt.
.3 long.txt.
.3 new1.sh.
.3 readonly.txt.
.3 script1.sh.
.3 sum.c.
.3 terminate.txt.
.2 dir2.
.3 csvfile1.csv.
.3 foo.txt.
.3 hello.txt.
.3 mysql.
.4 sql1.sql.
.2 files.txt.
.2 new.sh.
.2 old2.txt.
.2 recent1.txt.
.2 results.txt.
.2 test.
.3 2dir.
.3 3dir.
.3 dir1.
}
\caption{File System 3}
\end{subfigure}
\caption{File System structures designed for \benchmark{}-Bash.}
\label{fig:nl2b_filesystems}
\end{figure}

\textbf{Reward function.} Evaluation of an agent's trajectory across a single task episode towards carrying out the given instruction is determined by modifications to the file system and the latest execution output. The instructions found in the \benchmark{}-Bash dataset fall under one of two buckets: it either 1. Requests information about the file system that can be answered via execution output generated from a correct sequence of Bash actions (i.e. "How many files...", "What is the size of...", "Where is the .png image stored?") or 2. Requests a change to the location, configuration, or content of a file or folder (i.e. "Move the \texttt{dir1} folder from...", "Set the permissions to...", "Append a line to..."). Any relevant correct changes are therefore captured by considering both execution output and file system modifications during evaluation.

We define $A$ and $G$ as the outputs of the \texttt{agent} and \texttt{gold} commands respectively, where $A_{out}$ and $G_{out}$ refer to the execution output, and $A_{fs}$ and $G_{fs}$ refer to a list of entries reflecting file system modifications, where each entry is \texttt{[file path, modification type $\in$ [added, changed, deleted]]}. We then formally define the reward function as follows:

\begin{equation}
    \label{eq:reward_bash}
    \begin{gathered}
        \mathcal{R} = 0.34 * \texttt{similarity}(A_{out}, G_{out}) \\
        + 0.33 * (1 - \texttt{erf}(\lvert A_{fs} \cup G_{fs} - A_{fs} \cap G_{fs} \rvert)) + \\
        + 0.33 * \frac{\texttt{is\_correct}(A_{fs} \cap G_{fs})}{A_{fs} \cap G_{fs}}
    \end{gathered}
\end{equation}

Where \texttt{similarity} refers to lexical similarity, which is determined by the cosine similarity score between TF-IDF vectors (calculated with \texttt{TfidfVectorizer} from \texttt{scikit-learn}) of the two execution outputs. The second component of the reward function reflects the number of file system modifications that were either not completed or not necessary; the error associated with the total number of misses is constrained to the range \texttt{[0,1]} using the Gauss error function (\texttt{erf}), where $0$ corresponds to no file system modification mistakes. The third component checks what proportion of paths altered by both \texttt{agent} and \texttt{gold} were modified correctly. The \texttt{is\_correct} function returns the number of file paths that were changed correctly, determined by checking whether the \texttt{md5sum} hashes of each file path are identical for \texttt{agent} and \texttt{gold}. If $A_{fs} \cap G_{fs} = \emptyset$, this reward is automatically $1$. The scalar weights for each component are arbitrarily assigned.

A max score of $1$ is achieved only if the correct file paths are changed, the changes are correct, and the latest execution output matches the \texttt{gold} command output exactly. Figure~\ref{fig:preview} visualizes the reward function.
While an exact match comparison would have been a simpler choice to satisfy the Success Rate metric put forth in the main paper, we design this reward function to 1. Demonstrate that \benchmark{} can support complex reward functions that account for multiple forms of execution output, and 2. Provide practitioners who use the \benchmark{}-Bash environment with a scalar reward that reflects how "similar" the given output is to the expected, rather than a flat 0/1 reward value that may over-penalize and discount the efforts of more capable reasoning abilities. These reasons also motivate the SQL-based environment's reward function, discussed in the following section.

\begin{figure}
    \centering
    \begin{minipage}{0.49\textwidth}
        \centering
        \includegraphics[width=\textwidth]{figures/assets/bash_util_freq.pdf}
        \caption{Top 30 most frequently occurring bash utilities out of the 66 in \benchmark{}-Bash with their frequencies in log scale. }
        \label{fig:nl2b_utility_distribution}
    \end{minipage}\hfill
    \begin{minipage}{0.49\textwidth}
        \centering
        \includegraphics[width=\textwidth]{figures/assets/sql_hardness.png}
        \caption{Distribution of \texttt{gold} command difficult for \benchmark{}-SQL task data adapted from the Spider SQL dataset.}
        \label{fig:sql_hardness_distribution}
    \end{minipage}
\end{figure}

\subsection{SQL Environment}
\label{appx:env:sqlenv}
\textbf{Environment Definition.} The \texttt{Dockerfile} defining the SQL-based environment inherits from the MySQL image and adds a \texttt{.sql} file setup script to the \texttt{/docker-entrypoint-initdb.d} directory within the Docker image; this is a special directory made for container initialization. On container start-up, the added \texttt{.sql} file, which creates and populates databases with tables and tables with records, is automatically invoked. Since the \benchmark{}-SQL dataset does not feature any queries that involve modifying the database in any manner (i.e. no \texttt{INSERT}, \texttt{UPDATE}, or \texttt{DELETE} commands), there is no reset mechanism written into the \texttt{Dockerfile} definition that is invoked before each task episode; with that said, adding a reset script or version control to the \texttt{Dockerfile} is simple.

\textbf{\benchmark-SQL dataset.}
\benchmark-SQL is adopted from the development set of the Spider dataset \cite{yu-etal-2018-spider}. Spider 1.0 is a large-scale cross-domain dataset on generating SQL queries from natural language questions whose development set contains 1034 pairs of <\texttt{instruction}, \texttt{gold}> task instances spanning 20 databases.
The distribution of queries according to their hardness criterion is shown in Figure~\ref{fig:sql_hardness_distribution}.
As discussed in Section~\ref{sec:benchmark:ic_envs}, a filtering criterion narrows down the Spider dataset's information to only the necessary components.
We do not add anything to the Spider dataset that was not originally available.
The Spider 1.0 dataset is available for use under the \hyperlink{https://creativecommons.org/licenses/by-sa/4.0/legalcode}{CC BY-SA 4.0 license}.

\begin{figure}
    \centering
    \begin{minipage}{0.48\textwidth}
        \centering
        \includegraphics[width=\textwidth]{figures/assets/Traj-SQL}
        \caption{Example of interactions between an agent and the \benchmark{} SQL Environment}
        \label{fig:trajectory_sql}
    \end{minipage}\hfill
    \begin{minipage}{0.48\textwidth}
        \centering
        \includegraphics[width=\textwidth]{figures/assets/Reward-SQL}
        \caption{Evaluation of the results of agent interactions with the SQL Environment against the gold command associated with the task. A simple Intersection over Union formula that accounts for duplicates is used to quantify answer correctness. Task completion is a reward of 1.}
        \label{fig:reward_sql}
    \end{minipage}
\end{figure}

\textbf{MySQL databases.} We first resolve data types for primary, foreign key pairs across the provided table schemas in Spider for conflicting instances and generate the corresponding SQLite databases. Next, to align with our Docker-supported environment, we convert the SQLite databases to MySQL format using \texttt{sqlite3mysql}~\cite{sqlite3mysql}, a Python library, and then generate a unified MySQL dump having schemas for all the tables. To handle case-sensitive table name discrepancies between the queries and the underlying schema in the original Spider dataset, we activate the \texttt{lower\_case\_table\_names} setting in our evaluation environment. Additionally, for proper access controls, we create a test user and grant them all privileges for all the tables.

\textbf{Reward function.} The completion evaluation mechanism compares the output of the \texttt{gold} SQL command with the latest execution output (i.e. latest observation) from the agent's interaction trajectory. The execution output of all \texttt{gold} SQL queries is a list of records. Each record is a tuple of one or more values that may be different types. For any single execution output, the order of types for every record is identical. Given the \texttt{agent} command(s)' latest execution output $A$ and the \texttt{gold} command's execution output $G$, we formulate the reward function as follows:

\begin{equation}
    \label{eq:reward_sql}
    \mathcal{R} = \frac{A \cap G}{A \cup G} * (kendalltau((A \cap (A \cap G)), (G \cap (A \cap G))) + 1) / 2
\end{equation}

We employ Intersection over Union (\textit{IoU}), or more formally the Jaccard Index, to quantify the correctness of the latest execution output generated by the agent against the \texttt{gold} output. If the latest execution output of the SQL query is not in the form of a list of records (i.e. a string error message), the reward is $0$ by default. Among the items that lie in the intersection of the \texttt{agent} and \texttt{gold} execution outputs, we also apply a penalty if the records are in the incorrect order. Since achieving the correct order of fields in a record is of non-trivial importance to addressing many SQL queries correctly, we do not do any re-ordering or pre-processing of the list of records. Therefore, a record formatted as \texttt{("Ross", 29)} is not awarded any credit against a \texttt{gold} output that includes \texttt{(29, "Ross")}. To quantify how sorted the \texttt{agent} output is relative to the \texttt{gold} output, we lean on Kendall's $\tau$ and adjust the output range to $[0, 1]$. The \textit{IoU} score is then directly scaled by this coefficient.

All in all, only a correctly ordered list with the exact set of records found in the \texttt{gold} output would receive a max score of 1, which corresponds to task completion.
Figure~\ref{fig:reward_sql} visualizes the reward function for an example set of outputs. Note that in the main paper, the Success Rate metric is used; the scalar $3/7$ output shown in the figure is treated as a 0 when quantifying whether the task was completed via the 0/1 Success Rate metric. As mentioned in the discussion of the Bash reward function, this reward function also aims to be a richer and fairer continuous evaluation metric of a model's reasoning abilities compared to a binary 0/1 task completion score.

\subsection{Python Environment}
\textbf{Environment definition.} The \benchmark{}-Python task environment inherits from a bare-minimum Python 3.9 image that provides the basic essentials for initializing a Python interpreter.
We were unable to determine how to initialize a Python interpreter within a Dockerfile such that the container would then be capable of automatically executing Python commands sent to it while continuous logging every action/observation per turn.
To overcome this, we create and define a backend application that runs within the Docker container, simulates a Python interpreter, and is responsible for handling input/output.
By having the application sit between the agent's actions and the interpreter, we are able to log every episode faithfully in addition to providing an environment that is agent-friendly and faithful to the experience of a real Python interpreter.

\textbf{\benchmark{}-Python dataset.} A large majority of code datasets popular within the NLP community are based on Python and present code completion as the primary task~\cite{chen2021evaluating,austin2021program,hendrycksapps2021}.
In the original problem setting, a task worker is asked to synthesize code in a zero, one, or few-shot setting with little to no access to an execution environment.
In the interactive setting, task workers are asked to accomplish the same objective, but informed that they have a Python environment to do whatever may help them write the function correctly, such as prototype different implementations and write/execute their own unit tests.
Therefore, datasets such as HumanEval, APPS, and MBPP require little to no revisions to be usable within the \benchmark{} environment, with the only necessary processing for all three being renaming of dataset attributes to \benchmark{}-compatible names.
A visualization of an example trajectory of interactions between an agent and the Python interpreter is presented in Figure~\ref{fig:trajectory_python}.

\textbf{Reward function.} We preserve the original metric of proportion of unit tests passed to evaluate agent implementations, with all tests passing being equivalent to task completion.
Complementary to the visualization of interactions, we also show how \benchmark{}-Python performs automatic evaluation of an agent's implementation of the desired function in Figure~\ref{fig:reward_python}.

\begin{figure}
    \centering
    \begin{minipage}{0.48\textwidth}
        \centering
        \includegraphics[width=\textwidth]{figures/assets/Traj-Py.pdf}
        \caption{Example of interactions between an agent and the \benchmark{} Python Environment}
        \label{fig:trajectory_python}
    \end{minipage}\hfill
    \begin{minipage}{0.48\textwidth}
        \centering
        \includegraphics[width=\textwidth]{figures/assets/Reward-Py.pdf}
        \caption{In this setting, an agent interacts with a Python Interpreter to 1. implement the requested method and 2. write test cases to determine function correctness. Upon submission, the reward function then evaluates the agent's implementation with a set of unit tests.}
        \label{fig:reward_python}
    \end{minipage}
\end{figure}
\section{Experiment Details}
\label{appx:experiments}

\subsection{Model Details}
\label{appx:experiments:details}

We do not perform any model training for configuring the methods or running the experiments discussed in this project. Our evaluations use inference call requests to OpenAI, PaLM, and HuggingFace API endpoints to run the baseline models on the \benchmark{} tasks. For OpenAI models, we set temperature to $0$, \texttt{top\_p} to $1$, \texttt{max\_tokens} to 512, and \texttt{n} (number of completions) to $1$. For PaLM models, we set temperature to $0$, \texttt{top\_p} to $1$, and \texttt{candidate\_count} (number of completions) to $1$. For open source models, we set \texttt{max\_new\_tokens} (maximum number of tokens to generate) to $100$ and temperature to $0.01$. Due to constraints in the context window size, we limit the length of each observation to a maximum of $1000$ tokens across all inference calls. The code for configuring API calls can be found in the linked repository.

\subsection{Additional Experiments \& Analysis}
\label{appx:experiments:additional_expr}

% \textbf{\benchmark{}-SQL additional results.}
% Table~\ref{table:single-vs-multi-sql-extra} includes results for additional experiments on open-source models that were completed by the supplementary deadline. Again, the advantage of interactions offered in the Try Again scenario accounts for a multi-point disparity in successful task completion for the SQL task across all degrees of difficulty.

% \begin{table}[H]
\centering
\begin{tabular}{llllll|lllll}
\toprule
      \benchmark{}-SQL & \multicolumn{5}{c}{Single Turn} & \multicolumn{5}{c}{Try Again ($n$ = 10)} \\
     Model / Hardness & Easy & Med & Hard & Extra & All & Easy & Med & Hard & Extra & All \\
\midrule
Vicuna-13B & 8.1 & 1.3 & 0.6 & 0.0 & 2.6 & 18.9 & 3.4 & 1.7 & 0.0 & 6.3 \\
StarChat-16B & 21.8 & 7.4 & 2.9 & 0.0 & 8.9 & 22.3 & 8.5 & 2.9 & 1.2 & 9.7 \\
\bottomrule
\end{tabular}
\vspace{0.5em}
\caption{(Additional Results) Success Rate for single vs.\,multi turn evaluation on \benchmark{}-SQL (refer \S\ref{appx:env:sqlenv}). Query difficulty is adopted from Spider~\cite{yu-etal-2018-spider}.}
\label{table:single-vs-multi-sql-extra}
\vspace{-1.5em}
\end{table}

\textbf{SQL schema ablation.} To confirm that the benefits of interaction exceed a simple disparity in information between the Single Turn and Try Again settings, we add the full SQL database schema, providing holistic details of tables necessary to the given instruction, to the \texttt{Question} message of both prompts, then re-run the comparison for several. Table~\ref{table:ablation-handicap-sql} indicates that while Single Turn performance improves drastically, a non-trivial difference in favor of Try Again remains. Manual inspection of task episode trajectories shows that selective and fine-grained context discovery (i.e. inspecting specific table records and file content that affect query construction) is still critical to solving tasks efficiently.

\begin{table}[H]
\centering
\begin{tabular}{llllll|lllll}
\toprule
      \benchmark-SQL \\ \quad + Schema &
      \multicolumn{5}{c}{Single Turn} &
      \multicolumn{5}{c}{Try Again (max 10 turns)} \\
     Model / Hardness & Easy & Med & Hard & Extra & All & Easy & Med & Hard & Extra & All\\
\midrule
gpt-3.5-turbo & 90.7 & 70.2 & 59.2 & 37.3 & 67.9 & \textbf{92.7} & \textbf{74.9} & \textbf{67.2} & \textbf{43.4} & \textbf{72.8} \\
text-bison-001 & 89.5 & 68.2 & 44.2 & 19.3 & 61.4 & 90.7 & 70.4 & 50.0 & 21.1 & 63.9 \\
chat-bison-001 & 79.0 & 52.0 & 32.1 & 15.1 & 49.2 & 82.2 & 56.0 & 42.5 & 24.1 & 54.9 \\
\bottomrule
\end{tabular}
\vspace{0.5em}
\caption{Success Rate across difficulty for single vs. multi-turn evaluation on the \benchmark{}-SQL dataset, with the database schema relevant to each task episode's instruction, also provided in the \texttt{Question} message of the prompting strategy. Best metrics are in \textbf{bold}.}
\label{table:ablation-handicap-sql}
\vspace{-1.5em}
\end{table}

\textbf{Trends of admissible actions.}
Table~\ref{table:error-rate} shows that for the SQL task, models generate admissible actions with increasingly higher rates early on; in initial turns, models will tend to hallucinate a query with fabricated table and column names at a high frequency. The drop in error rate between the first and second turns can largely be attributed to the model's decision to begin exploring context; $60.3\%$ of second turn actions contain either the \texttt{SHOW TABLES} or \texttt{DESC} keywords. Prompting strategies (i.e. ReAct, Plan \& Solve), explicit phrasing that encourages exploration, and demonstrations diminish a model's default tendency to hallucinate a query in the first turn. This trend is not found in Bash. This can likely be attributed to the nature of the instructions; unlike the SQL instructions which simply pose a question and do not have any explicit references to SQL commands or clauses, Bash instructions will typically include keywords that correspond directly to useful Linux commands or give insight into the file system's internal structure. These signals reduce the need for context discovery. Therefore, successful task completion in Bash tends to lean towards 1) Figuring out which flags, options, and arguments to configure a command with and 2) How to string together commands or pass outputs from one command to the next correctly.

For both Bash and SQL, in later turns, the rate of admissible actions does not improve consistently. The actions in these later turns are usually attempts to answer the original instruction.
At these stages, a model will tend to make small, cursory adjustments to the prior action based on execution feedback, often resulting in both a repetition of the same types of mistakes and hallucinations that introduce new issues.
In these moments, compared to such minor perturbations, alternative reasoning capabilities such as context discovery and modularized problem solving are often more efficient ways to get the relevant insights needed to better decide how to fix the prior turns' issues.
As corroborated by Figure~\ref{fig:turn_stats}, models struggle to take advantage of additional context in longer task episodes or horizons.
Making the most of multiple queries is an open challenge with exciting implications for solving more difficult coding tasks.

\begin{table}[H]
    \centering
    \begin{tabular}{llllllllllll}
        \toprule
        Turn & 1 & 2 & 3 & 4 & 5 & 6 & 7 & 8 & 9 & 10 \\
        \midrule
        SQL  & 90.2 & 46.4 & 34.4 & 39.7 & 31.1 & 42.9 & 51.5 & 47.4 & 48.4 & 46.6 \\
        Bash & 23.1 & 28.6 & 34.7 & 37.5 & 37.6 & 42.9 & 39.3 & 37.1 & 33.7 & 38.2 \\
        \bottomrule
    \end{tabular}
    \vspace{0.5em}
    \caption{Error \% (Average ratio of non-admissible actions) per turn for the Try Again prompting scheme using a GPT 3.5 model on the Bash and SQL \benchmark{} datasets.}
    \vspace{-1.5em}
    \label{table:error-rate}
\end{table}

\textbf{Robustness results.}
We conducted an evaluation to assess the robustness of the reported accuracy metrics for the models. In order to maintain consistency in the evaluation, we focused on the performance across file systems 2, 3, and 4 (shown in Figure~\ref{fig:nl2b_filesystems}), which were designed to have similar difficulty levels. File system 1, intentionally made harder, was not included in this analysis.
The standard errors for the Single Turn and Try Again modes are presented in Table~\ref{table:bash-error-bars}. The Try Again mode leverages interaction to consistently outperform the Single Turn mode across all models.

\begin{table}[H]
\centering
\begin{tabular}{ll|ll}
\toprule
      Model & {Single Turn} & {Try Again ($n$ = 10)} \\
\midrule
    text-davinci-003 & $31.40 \pm 1.35$ & $43.13 \pm 5.98$ \\
    gpt-3.5-turbo & $36.63 \pm 1.83$ & $47.40 \pm 1.23$\\
    gpt-4 & $\mathbf{38.37 \pm 1.20}$ & $\mathbf{52.70 \pm 3.50}$\\
    text-bison-001 & $18.83 \pm 3.57$& $22.40 \pm 3.35$\\
    chat-bison-001 & $20.47 \pm 1.89$& $21.67 \pm 1.81$\\
    Vicuna-13B & $16.73 \pm 5.00$ & $27.67 \pm 4.15$\\
    StarChat-16B & $19.37 \pm 3.04$ & $27.17 \pm 2.74$\\
\bottomrule
\end{tabular}
\vspace{0.5em}
\caption{(Robustness Results) Success Rate with standard errors for single vs.\,multi turn evaluation on \benchmark{}-Bash (refer \S\ref{appx:env:bashenv}). Best metrics are in \textbf{bold}. Both modes display significant standard errors (as expected) but still Try Again outperforms Single Turn by a huge margin.
}
\label{table:bash-error-bars}
\vspace{-1.5em}
\end{table}

\subsection{Additional Prompting Strategy}
\label{appx:experiments:additional_strategies}

To gauge the significance of designing prompting strategies that can successfully solve the interactive coding task, we attempt to devise a more performant approach by chaining together existing techniques, where each technique is meant to elicit a different, relevant reasoning skill. To this end, we design a hybrid prompting strategy that combines Plan \& Solve and Try Again, which we refer to as "Plan \& Solve + Refine". This strategy is meant to complement a model's planning, modularized task completion, and context discovery abilities with error correction. Figure~\ref{fig:psr_visual} visualizes this prompting strategy's workflow. The full prompting template is included in \S~\ref{appx:experiments:templates}.

\begin{figure}[H]
    \centering
    \includegraphics[width=1.0\textwidth]{figures/assets/PSR.pdf}
    \caption{Visualization of the Plan \& Solve + Refine strategy. This prompting strategy is a naive combination of Plan \& Solve and Try Again; the model is first asked to formulate and execute a plan. If it fails, then the model is asked to continue to work on generating the right solution until the reward is $1$ or the number of allowed turns ($n$ = 3) is exhausted.}
    \label{fig:psr_visual}
    \vspace{-1em}
\end{figure}

We run small-scale experiments in the same fashion as discussed in Section~\ref{sec:experiments:strategy_compare} on the SQL task for $44$ randomly selected task instances. These preliminary tests have an average Success Rate of 45.5, a Number of Turns of 6.34 (across both planning and refinement stages), and an Error \% of 22.4. Compared to the original Plan \& Solve performance presented in Section~\ref{sec:experiments:strategy_compare}, this strategy does not yield a significant improvement. Upon inspection of the trajectories, it is evident that while models are capable of specific reasoning techniques, it is not enough to elicit specific strategies in a pre-determined order. In addition, compared to human problem solvers, such reasoning frameworks do not intuitively allow for adjustment in strategy based on execution output. Therefore, mechanisms that allow a model to freely self-determine what strategies at any turn in task episode are integral to improving Success Rate and non-trivial because simple adjustments in existing prompting strategies are not enough to close this gap. 

\subsection{Task Trajectory Analyses}
\label{appx:experiments:trajectory_analyses}

In this section, we discuss some of the behaviors and phenomena in different prompting schemes. We also include annotated trajectories that are representative of these observations.

\begin{table}[H]
\footnotesize
\begin{minipage}[t]{.49\linewidth}
\begin{tabularx}{\textwidth}{X}
\toprule
\textbf{Inst}: Count md5sum of all '*.py' files in /testbed folder with subfolders. \\
\textbf{Gold}: \texttt{find /testbed -type f -name "*.py " -exec md5sum \{\} + \text{\textbar} awk '\{print \$1\}' \text{\textbar} sort \text{\textbar} md5sum} \\
\textbf{Out}: \texttt{ff48e395ab4e0acb8336c4afe4504  -} \\
\midrule
\textbf{Act}: \textcolor{red}{\texttt{find /testbed -name "*.py" -type f -exec md5sum \{\} + \text{\textbar} wc -l}} \\
\textbf{Obs}: 4 \\
\textbf{Act}: \texttt{find /testbed -name "*.py" -type f -exec md5sum \{\} + \text{\textbar} \textcolor{orange}{awk 'END\{print NR\}'}} \\
\textbf{Obs}: 4 \\
\textbf{Act}: \texttt{find /testbed -name "*.py" -type f -exec md5sum \{\} + \text{\textbar} \textcolor{orange}{awk 'END\{print\}' \text{\textbar} cut -d' ' -f1}}  \\
\textbf{Obs}: 0d1e201ac67b5a979c52f42eb8dd3702 \\
\textbf{Act}: \texttt{find /testbed -name "*.py" -type f -exec md5sum \{\} + \text{\textbar} \textcolor{orange}{awk '\{print \$1\}' \text{\textbar} sort \text{\textbar} md5sum \text{\textbar} cut -d' ' -f1}} \\
\textbf{Obs}: 7e9ff48e395ab4e0acb8336c4afe4504 \\
% \textbf{Act}: \texttt{find /testbed -name "*.py" -type f -exec md5sum \{\} + \text{\textbar} awk '\{print \$1\}' \text{\textbar} \textcolor{orange}{md5sum \text{\textbar} cut -d' ' -f1}} \\
% \textbf{Obs}: 3f46f161d2dea7e424b2f2f7f887ee0c \\
... (6 more turns) \\
\bottomrule
\end{tabularx}
\end{minipage}
\begin{minipage}[t]{.49\linewidth}
\begin{tabularx}{\textwidth}{X}
\toprule
\textbf{Inst}: For each continent, list its id, name, and how many countries it has? \\
\textbf{Gold}: \texttt{SELECT T1.ContId , T1.Continent , count(*) FROM CONTINENTS AS T1 JOIN COUNTRIES AS T2 ON T1.ContId = T2.Continent GROUP BY T1.ContId;} \\
\midrule
\textbf{Act}: \textcolor{red}{\texttt{SELECT continent.id, continent.na- me, COUNT(country.id) AS num\_countries FROM continent LEFT 
JOIN country ON continent.id = country.continent\_id GROUP BY continent.id, continent.name}} \\
\textbf{Obs}: Error: Table 'car\_1.continent' doesn't exist \\
\textbf{Act}: \textcolor{blue}{\texttt{SHOW TABLES}} \\
\textbf{Obs}: [('car\_makers',), ('car\_names',)] + 4 more... \\
\textbf{Act}: \textcolor{blue}{\texttt{DESC continents}} \\
\textbf{Obs}: [('ContId', 'int', 'NO', 'PRI', None, 'auto\_increment'), ('Continent', 'text', 'YES'...
\\
\textbf{Act}: \texttt{SELECT \textcolor{orange}{ContId, Continent, COUNT(*) AS NumCountries} FROM \textcolor{orange}{countries} GROUP BY \textcolor{orange}{Continent}} \\
% \textbf{Obs}: Error: Unknown column 'ContId' \\
% \textbf{Act}: \texttt{SELECT \textcolor{red}{Id}, Continent, COUNT(*) AS NumCountries FROM countries GROUP BY Continent} \\
...(5 more turns) \\
\bottomrule
\end{tabularx}
\end{minipage}
\vspace{1em}
\caption{Examples of Try Again interaction trajectories with GPT 3.5 based model from Bash (left) and SQL (right). Without any reasoning framework, models will tend to \textcolor{red}{hallucinate} an initial command, then make incremental adjustments (\textcolor{orange}{error correction}) to the command based on the execution feedback. On occasion, a model will write commands to \textcolor{blue}{gather context} about the task environment. For both trajectories, $n$ = 10 turns was allotted and reward was $0$ (max turns exhausted).}
\label{table:expr-default-behavior}
\vspace{-1.5em}
\end{table}

\textbf{Default model behaviors.} In both the Single Turn and Try Again setups, the templates describe the interactive coding task and do not provide the model with any reasoning framework or strategy, providing insight into a model's default approach when given the instruction and task environment. As evident in Table~\ref{table:expr-default-behavior}, across both the Bash and SQL tasks, a model tends to directly generate an answer in the very first turn. For Try Again, the model will then perform error correction, then make slight adjustments to the answer to resolve the issue(s) reflected in prior execution output across the next several turns. With less frequency, for SQL, the model will also issue commands to get database and table descriptions. To the best of our knowledge, based on a close examination of the Try Again trajectories, current language models do not demonstrate any innate reasoning capacity for the interactive coding task beyond error correction and context discovery.

\begin{table}
\footnotesize
\textbf{Inst}: What is TV Channel of TV series with Episode "A Love of a Lifetime"? List TV Channel's series name. \\
\textbf{Gold}: \texttt{SELECT T1.series\_name FROM TV\_Channel AS T1 JOIN TV\_series AS T2 ON T2.Channel = T1.id WHERE T2.Episode = "A Love of a Lifetime";} \\
\\
% \scriptsize
\begin{minipage}[t]{.49\linewidth}
\begin{tabularx}{\textwidth}{X}
\toprule
\textbf{Try Again} (r = 0, length = 10, max turns = 10) \\
\midrule
\textbf{Act}: \texttt{\textcolor{violet}{SELECT t1.channel, t2.series\_name FROM \textcolor{red}{table1} t1 JOIN \textcolor{red}{table2} t2 ON t1.\textcolor{red}{series\_id} = 
t2.\textcolor{red}{series\_id} WHERE t2.\textcolor{red}{episode\_name} = 'A Love of a Lifetime'}} \\
\textbf{Obs}: Error: Table 'table1' doesn't exist \\
% \vspace{0.1pt}
\textbf{Act}: \textcolor{orange}{\texttt{SHOW TABLES}} \\
\textbf{Obs}: [('cartoon',), ('tv\_channel',), ('tv\_series',)] \\
% \vspace{0.1pt}
\textbf{Act}: \textcolor{orange}{\texttt{DESC tv\_channel}} \\
\textbf{Obs}: [('id', 'varchar(255)', 'NO', 'PRI', None, ''), ('series\_name', 'text', 'YES', '', None...] + 7 
more... \\
% \vspace{0.1pt}
\textbf{Act}: \textcolor{orange}{\texttt{DESC tv\_series}} \\
\textbf{Obs}: [('id', 'double', 'NO', 'PRI', None, ''), ('Episode', 'text', 'YES', '' + 7 more... \\
% \vspace{0.1pt}
\textbf{Act}: \texttt{\textcolor{violet}{SELECT \textcolor{red}{Channel},series\_name FROM tv \_channel WHERE id IN (SELECT Channel FROM tv\_series 
WHERE Episode = 'A Love of a Lifetime')}} \\
\textbf{Obs}: Error: Unknown column 'Channel' \\
% \vspace{0.1pt}
\textbf{Act}: \textcolor{violet}{\texttt{SELECT \textcolor{red}{channel\_name},series\_name FROM tv\_channel JOIN tv\_series ON tv\_channel.id = 
tv\_series. Channel WHERE Episode = 'A Love of a Lifetime'}} \\
\textbf{Obs}: Error: Unknown column 'channel\_name' \\
% \vspace{0.1pt}
\textbf{Act}: \textcolor{violet}{\texttt{SELECT Channel,series\_name FROM tv\_channel JOIN tv\_series ON tv\_channel.id = 
tv\_series.Channel WHERE Episode = 'A Love of a Lifetime'}} \\
% \textbf{Obs}: [('700', 'Sky Radio')] \\
% % \vspace{0.1pt}
% \textbf{Act}: \textcolor{violet}{\texttt{SELECT tv\_channel.channel\_name, tv\_series.se-
% ries\_name FROM tv\_channel JOIN tv\_series ON tv\_c-
% hannel.id = tv\_series.Channel  WHERE Episode = 'A Love of a Lifetime'}} \\
% \textbf{Obs}: Error: Unknown column 'channel\_name' \\
% \vspace{0.1pt}
Continued hallucinations in problem solving actions for last 4 turns... \\
\bottomrule
\end{tabularx}
\end{minipage}
\begin{minipage}[t]{.49\linewidth}
\begin{tabularx}{\textwidth}{X}
\toprule
\textbf{ReAct} (r = 1, length = 7, max turns = 10) \\
\midrule
\textbf{Thought}: I should check what tables are available in this MySQL database that can help me answer this 
question \\
\textbf{Act}: \textcolor{blue}{\texttt{SHOW TABLES}} \\
\textbf{Obs}: [('cartoon',), ('tv\_channel',), ('tv\_series',)] \\
% \vspace{0.1pt}
% \textbf{Thought}: I should check if the tv\_series table might have any information related to the episode name or 
% the TV Channel \\
% \textbf{Act}: \textcolor{blue}{\texttt{DESC tv\_series}} \\
% \textbf{Obs}: [('id', 'double', 'NO', 'PRI', None, ''), ('Episode', 'text', 'YES', '', None, '')] + 7 more... \\
% % \vspace{0.1pt}
...(1 turn for context discovery with \texttt{DESC}) \\
% \vspace{0.1pt}
\textbf{Thought}: The tv\_series table has a Channel column that might be helpful. I should check if there is a row with the episode name "A Love of a Lifetime". \\
\textbf{Act}: \textcolor{violet}{\texttt{SELECT Channel FROM tv\_series WHERE Episode = 'A Love of a Lifetime'}} \\
\textbf{Obs}: [('700',)] \\
% \vspace{0.1pt}
\textbf{Thought}: I get a single value, 700, which is the TV Channel of the TV series with the episode "A Love of a
Lifetime". Check what series is on this channel. \\
\textbf{Act}: \textcolor{violet}{\texttt{SELECT series\_name FROM tv\_channel WHERE \textcolor{red}{channel\_number} = 700}} \\
\textbf{Obs}: Error: Unknown column 'channel\_number' \\
% \vspace{0.1pt}
\textbf{Thought}: I should check what columns are available in the tv\_channel table. \\
\textbf{Act}: \textcolor{orange}{\texttt{DESC tv\_channel}} \\
\textbf{Obs}: [('id', 'varchar(255)', 'NO', 'PRI', None, ''), ('series\_name', 'text', 'YES', ''... + 7 more... \\
% \vspace{0.1pt}
\textbf{Thought}: The tv\_channel table does not have a channel\_number column. Check if there is another way to get the series name for the TV Channel 700. \\
\textbf{Act}: \textcolor{violet}{\texttt{SELECT series\_name FROM tv\_channel WHERE id = '700'}} \\
\textbf{Obs}: [('Sky Radio',)] \\
% \vspace{0.1pt}
\textbf{Thought}: I get a single value, Sky Radio, which is the series name for the TV Channel 700. \\
\textbf{Act}: \textcolor{violet}{\texttt{submit}} \\
\bottomrule
\end{tabularx}
\end{minipage}

\vspace{0.5em}
\caption{Comparison of Try Again \& ReAct interaction trajectories with GPT 3.5 based model towards solving the same instruction. Without explicit strategy or task demonstrations for Try Again, the model tends to \textcolor{red}{hallucinate} and doesn't build upon prior observations as well. With ReAct, there is greater evidence of \textcolor{blue}{context discovery}, \textcolor{orange}{error correction}, and \textcolor{violet}{problem solving}.}
\label{table:expr-traj-analysis}
\vspace{-2em}
\end{table}

\textbf{Prompting strategy effects.} In contrast with Try Again, the ReAct prompting strategy briefly introduces the interactive task environment and proposes a reasoning framework for the model to abide by. Table~\ref{table:expr-traj-analysis} presents a side-by-side comparison of the Try Again and ReAct~\cite{yao2023react} strategies. The figure reflects the richer types of problem-solving that a model exhibits when prompted to reason on its own thoughts about the execution output. This reflects that through better strategies, it may be possible to make significant advancements in the interactive coding task with prompting strategies that attempt to elicit reasoning via an appropriate framework that also permits the model to be expressive and creative in devising its own solutions. This is particularly necessary for interactive code tasks, which pose multiple challenges that cannot be overcome by any isolated reasoning technique. As demonstrated in \S~\ref{appx:experiments:additional_strategies}, this direction is non-trivial, and \benchmark{} is designed to facilitate the bench-marking of such approaches.

\subsection{Capture the Flag Analysis}
\label{appx:experiments:ctf_analysis}

CTF challenges typically necessitate a trial-and-error methodology, where participants employ diverse techniques and exploit vectors to identify vulnerabilities to solve challenges. Processes such as exploring complex environments or executables, debugging, and dynamic exploitation, which involve sequential steps, require iterative interaction. Considering the inherently interactive nature of the task, it is crucial for an agent to employ an iterative approach and have access to an interactive platform to achieve success. In most instances, both humans and agents find it impracticable to solve a challenge in a single attempt.

Here, we present a more thorough discussion of Figure~\ref{fig:traj_ctf}. It is important to note that without the provided hint regarding the usefulness of the "sleuthkit" library, the agent fails to solve the task and engages in incorrect reasoning. However, upon receiving the prompt's hint, the agent adeptly utilizes this information to install the library and leverage its functionalities for its advantage. By analyzing a given disk image file, the agent employs the "mmls" command to inspect the corresponding partition table. From the partition table, it deduces that a significant portion of the space remains unallocated, while a Linux partition initiates at sector 2048. Subsequently, the agent attempts to access the contents of this sector using the "fls" command, searching for the "down-at-the-bottom.txt" file, which it anticipates will contain the flag. When unable to locate the file, the agent speculates that a recursive search might be necessary and adds the "-r" flag to its command. Due to the immense output, it becomes arduous to track the file's location, prompting the agent to employ the "grep" command to search for the file within the output. By examining the grep output, the agent identifies the file's location (18291) and proceeds to inspect its contents. The flag, presented in a visual format, is accurately recognized and submitted by the agent.

A human expert employs a very similar approach when provided with the hint. By furnishing an interactive framework, \benchmark{} empowers agents to emulate human-like behavior, enabling them to explore the environment, decompose tasks into subtasks, debug using traces and logs, and iteratively accumulate knowledge to successfully solve challenges.

\subsection{Human Performance Baseline}
\label{appx:experiments:human_perf}

To explore the gap between human and agent performance on the interactive coding task, we the authors, all proficient in SQL, act as human task workers and perform the task on a random sample of $15$ \benchmark{}-SQL task instances within the same task environment identical to the agent's setting. A max number of $n$ = $10$ turns is imposed, as was done with the Try Again prompting strategy. Similar to ReAct and Plan \& Solve, the human task worker decides when to \texttt{submit}; in other words, the task does not terminate automatically when reward = $1$. The trajectories for these $15$ instances and the code for facilitating human interaction with the \benchmark{}-SQL environment are available in the codebase.

The human task worker was able to complete 13 of 15 tasks (Success Rate = $0.87$) with low Error \%, most of the errors occurring not because of hallucinations of table columns and attributes, but rather because of SQL syntax errors that arose due to mistakes in relatively complex queries. What's noteworthy about the human task worker's trajectories is the presence of much more modularized problem-solving that deviates heavily from an agent's approach of generating a query in a single go. Even with context discovery and error correction, an agent's action to produce an answer for the instruction will tend to be a single, self-contained command that generates the answer in one go. On the other hand, a human task worker will tend to break up the query solution into multiple smaller sub-problems. This is particularly evident for instructions that must be answered with investigations across multiple tables with relations established by primary and foreign key columns. As an example, given an instruction "Find the average weight of the dog breed that is owned by the majority of pet owners", a human task worker might write commands that query the \texttt{pet\_owners} table to determine what the most popular dog breed is, and then use the answer to this sub-problem as a field in the \texttt{WHERE} clause of a second query that then determines the average weight using the \texttt{pets} table.

A more thorough and variegated study would be required to fully establish the performance gap between humans and agents. Nevertheless, from this small study, we are confident that humans generally exhibit more flexible and variegated reasoning capabilities compared to agents in the interactive coding task. Closing this gap is an exciting research direction, and beyond model-side improvements and scaling laws, incorporating human task reasoning and execution as guidance, feedback, or reward signals is a worthwhile consideration toward improving model performance.

\subsection{Prompt Templates}
\label{appx:experiments:templates}
As discussed in the paper, the main baseline evaluations for \benchmark{} consist of presenting a language agent with an instruction and a prompting strategy that have been adapted for \benchmark{}'s interactive task setting. Each prompting strategy is defined as a template with three components:

\begin{itemize}[noitemsep,topsep=0pt]
    \item Initial Message: This is the first message presented to the agent. The initial message may describe the general task to accomplish, guidelines for interacting with the \benchmark{} environment, the formats of the instruction and observation(s), and any additional information that pertains to the environment. In addition to the environment and task specifications, the general prompting strategy and useful demonstrations may also be discussed. The initial message is presented once as the first message of a task episode.
    \item Instruction Message: This is the template for communicating the instructions that an agent is asked to solve for a particular task episode. The instruction message is presented once as the second message of a task episode.
    \item Observation Message: This template is for communicating the standard output and any additional information for a single interaction. This observation is what the agent will use to generate the next action. The observation message may be presented multiple times depending on how many interactions the task episode lasts for.
\end{itemize}

Figures~\ref{appx:table:multi-turn-prompt}, ~\ref{appx:table:prompt-react}, ~\ref{appx:table:prompt-ctf}, and ~\ref{appx:table:prompt-plan-solve} present the corresponding prompt templates for the Try Again, ReAct, and Plan \& Solve experiments, along with a specific version for the toy Capture the Flag task.

\subsection{Supported Datasets}
While evaluation for Bash and SQL are carried out on the NL2Bash and Spider datasets, \benchmark{} supports multiple existing datasets based on these two languages and Python. We include Table~\ref{table:supported-datasets} to list all datasets currently supported by each \benchmark{} environment. Specific details regarding the transformation procedure and usage guidelines for each dataset can be found in the main \benchmark{} code repository.

\begin{table}[ht]
\centering
\begin{tabular}{l|l}
\toprule
    \benchmark{} Environment & Supported Datasets \\
\midrule
    IC-Bash & NL2Bash~\cite{lin-etal-2018-nl2bash} \\
    IC-Python & MBPP~\cite{austin2021program}, APPS~\cite{hendrycksapps2021} \\
    IC-SQL & Spider~\cite{yu-etal-2018-spider}, BIRD-SQL~\cite{li2023llm}, WikiSQL~\cite{zhong2017seq2sql} \\
\bottomrule
\end{tabular}
\vspace{0.5em}
\caption{Summary of all datasets supported by each \benchmark{} environment.}
\vspace{-1.5em}
\label{table:supported-datasets}
\end{table}

\section{Future Work Discussion}
In this section, we present some details of ongoing work to expand \benchmark{}'s coverage to more language, datasets, and tasks.

\textbf{Compiled language support.} Unlike interactive mode languages where expressions can be executed REPL-style one line at a time, imperative and interpreted languages that are typically processed by compilers (i.e. C, C++, Java, Go, Rust) are not as malleable to the exact form of the Bash or SQL environment.
To this end, we see two viable avenues of support for such languages:

\begin{itemize}[noitemsep,topsep=0pt]
    \item 3rd party interpreter support: Following Python, a language with both interpreter and compiler support, tools such as JShell (for Java) or Yaegi (for Go) may be serviceable interpreters for enabling REPL style code interaction for such languages. The main drawback to this approach is that this usage style feels a bit contrived and is not really found in real world software development processes.
    \item Multi-language environments: By creating an \benchmark{}-Bash based environment with a language’s corresponding compiler installed (i.e. \texttt{javac}, \texttt{gcc}), an agent would be able to use Bash commands to create, write to, and execute compiled-language files. (i.e. \texttt{touch hello.java; echo [cmd] > hello.java; javac hello.java; java hello}). While the execution of languages as an action in such a setting is not as direct as Option A, we believe that this paradigm is a practical setting that 1. Mirrors real world software engineering and 2. fits naturally with the interactive coding task formulation presented by \benchmark{}.
\end{itemize}

As a side note, Bash, Python, and SQL were the initial two languages chosen due to the bounty of such datasets that are already available.
On the contrary, despite their popularity among developers, there is a relative lack of such datasets for other languages (e.g., Java, C++, JavaScript) in the LLM2Code or NL2Code spaces. By 1. demonstrating interactive coding as a feasible, practical, and worthwhile task and 2. designing a language agnostic framework for task construction, we hope \benchmark{} might encourage more exploration into coding tasks that leverage interaction with 1+ programming languages that are not as popular at the moment.

\textbf{Beyond code generation.} It has been increasingly evident in recent years that many interactive tasks can be readily converted to Python-based code interaction problems, such as Python API interactions with a search engine to perform question answering or navigate websites for shopping~\cite{yao2022webshop}, code as interactive control policies for robots~\cite{liang2023code}, or code as a vehicle of thought for accomplishing complex, multi-step math problems~\cite{chen2023program}.
As code has become the medium of communication for many non-code synthesis tasks, we look forward to demonstrating and supporting \benchmark{}'s use as an environment for similar future tasks that extend into domains such as robotics, software engineering, and natural sciences.
\newcommand{\toprulec}{\arrayrulecolor{black}\specialrule{0.1em}{\abovetopsep}{0pt}%
            \arrayrulecolor{lightgray}\specialrule{\belowrulesep}{0pt}{0pt}%
            \arrayrulecolor{black}}
\newcommand{\midrulec}{\arrayrulecolor{lightgray}\specialrule{\aboverulesep}{0pt}{0pt}%
            \arrayrulecolor{black}\specialrule{\lightrulewidth}{0pt}{0pt}%
            \arrayrulecolor{white}\specialrule{\belowrulesep}{0pt}{0pt}%
            \arrayrulecolor{black}}
            
\begin{table}
    \centering
    \small
    \begin{tabular}{>{\raggedright\arraybackslash\tt}p{\textwidth}<{}}
        \toprulec
        \multicolumn{1}{c}{\cellcolor{lightgray} Initial Message}\\
        \midrulec
        \#\# TASK DESCRIPTION \\
        You are a {\color{blue} \{self.language\}} code generator helping a user answer a question using {\color{blue} \{self.language\}}.
        The user will ask you a question, and your task is to interact with a {\color{blue} \{self.setting\}} system using
        {\color{blue} \{self.language\}} commands to come up with the answer. \\
        \\
        \#\# RULES \\
        1. Do NOT ask questions \\
        2. Your response should only be {\color{blue} \{self.language\}} commands \\
        \\
        \#\# RESPONSE FORMAT \\
        Your response should be a {\color{blue} \{self.language\}} command. Format your {\color{blue} \{self.language\}} command as follows: \\
        \\
        \textasciigrave\textasciigrave\textasciigrave{\color{blue} \{self.language\}} \\
        Your {\color{blue} \{self.language\}} code here \\
        \textasciigrave\textasciigrave\textasciigrave \\
        \\
        Write {\color{blue} \{self.language\}} commands to help you do two things:
        \\
        1. Learn more about the {\color{blue} \{self.setting\}} you are interacting with. For example, if you are interacting with a MySQL database, you can use the DESCRIBE command to learn more about the tables
        you have access to. \\
        2. Execute {\color{blue} \{self.language\}} commands based on these tables to generate the correct output. \\
        \\
        DO NOT WRITE ANYTHING EXCEPT FOR CODE in your response. \\
        \\
        \#\# OUTPUT DESCRIPTION
        Given your {\color{blue} \{self.language\}} command input, the system will then give back output formatted as follows: \\
        \\
        Output: <string> \\
        Reward: [0, 1] \\
        \\
        The output is the standard output from executing your {\color{blue} \{self.language\}} command.
        The reward is a decimal value between 0 and 1, which tells you how close your {\color{blue} \{self.language\}}
        command is to the correct answer. The closer the reward is to 1, the closer your {\color{blue} \{self.language\}}
        command is to the correct answer.
        \vspace{0.5em} \\\toprulec
        \multicolumn{1}{c}{\cellcolor{lightgray} Instruction Message}\\
        \midrulec
        Query: "{\color{blue}<Task Episode Instruction>}"
        \vspace{0.5em} \\\toprulec
        \multicolumn{1}{c}{\cellcolor{lightgray} Observation Message}\\
        \midrulec
        Output: {\color{blue}<Observation from Intercode Environment>} \\
        Reward: {\color{blue}<Value $\in [0, 1]$>} \\
        \bottomrule
    \end{tabular}
    \vspace{3pt}
    \caption{Initial, Query, and Retry Prompts for the Try Again Strategy Evaluation. \texttt{\color{blue} \{self.language\}} is substituted with the programming language for interacting with the environment. \texttt{\color{blue} \{self.setting\}} is substituted with a natural language description of the environment}
    \label{appx:table:multi-turn-prompt}
\end{table}
            
\begin{table}
    \centering
    \small
    \begin{tabular}{>{\raggedright\arraybackslash\tt}p{\textwidth}<{}}
        \toprulec
        \multicolumn{1}{c}{\cellcolor{lightgray} Initial Message}\\
        \midrulec
        Interact with a {\color{blue}\{self.setting\}} system using {\color{blue}\{self.language\}} queries to answer a question.\\
        Use interleaving Thought, Action, Observation steps. \\
        Thought can reason about the current situation, and Action can be two types: \\
        (1) execute, which executes {\color{blue}\{self.language\}} code on the {\color{blue}\{self.setting\}} system \\
        (2) submit, which indicates that the previous observation is the answer \\
        Here are some examples. \\
        \\
        Question: What are the names and grades for each high schooler? \\
        Thought 1: I should write a SQL command that selects the name and grade fields from a table about high schoolers. \\
        Action 1: execute[SELECT name, grade FROM high\_schoolers] \\
        Observation 1: Error executing query: Table 'network\_1.high\_schoolers' doesn't exist \\
        Thought 2: The high\_schoolers table doesn't exist, I need to check what tables are available to use. \\
        Action 2: execute[SHOW TABLES] \\
        Observation 2: [('friend',), ('highschooler',), ('likes',)] \\
        Thought 3: There is a table called highschooler. I should check what information this table contains. \\
        Action 3: execute[DESC highschooler] \\
        Observation 3: [('ID', 'int', 'NO', 'PRI', None, 'auto\_increment'), ('name', 'text', 'YES', '', None, ''), 
        ('grade', 'int', 'YES', '', None, '')] \\
        Thought 4: The highschooler table contains the name and grade columns. I can now write a SQL command that answers the question \\
        ...
        \vspace{0.5em} \\\toprulec
        \multicolumn{1}{c}{\cellcolor{lightgray} Instruction Message}\\
        \midrulec
        Query: "{\color{blue}<Task Episode Instruction>}"
        \vspace{0.5em} \\\toprulec
        \multicolumn{1}{c}{\cellcolor{lightgray} Observation Message} \\
        \midrulec
        Output: {\color{blue}<Observation from Intercode Environment>} \\
        \bottomrule
    \end{tabular}
    \vspace{3pt}
    \caption{Initial, Query, and Retry Prompts for ReAct Strategy Evaluation. Additional ReAct style demonstrations written for SQL and Bash can be found on the \benchmark{} project GitHub.}
    \label{appx:table:prompt-react}
\end{table}

\begin{table}
    \centering
    \small
    \begin{tabular}{>{\raggedright\arraybackslash\tt}p{\textwidth}<{}}
        \toprulec
        \multicolumn{1}{c}{\cellcolor{lightgray} Initial Message} \\
        \midrulec
        You are an expert Capture the Flag (CTF) player.
        \vspace{0.5em} \\\toprulec
        \multicolumn{1}{c}{\cellcolor{lightgray} Instruction Message} \\
        \midrulec
        You are in the '/ctf' folder of a standard Ubuntu machine. {\color{blue} <Task Episode Instruction>}. Your goal is to find
this hidden flag by interacting with the bash terminal step-by-step. The flag is in the format 'picoCTF{{...}}'.
You can perform any valid action in the terminal by starting your action with 'cmd <your\_command>'. For bash
commands, prepend '/bin/bash -c \"<your\_command>\"', for python commands, prepend 'python3 -c \"<your\_command>\"'
and so on. For every action you take, you will receive the corresponding standard output as observation. You can
also explore/inspect the file system or do anything that helps you find the flag. Once you feel you have found
the flag, you can choose to submit it by printing 'submit <your\_flag>'. Do NOT provide an explanation for your
answer, only output the action you want.
        \vspace{0.5em} \\\toprulec
        \multicolumn{1}{c}{\cellcolor{lightgray} Observation Message} \\
        \midrulec
        Output: {\color{blue}<Observation from Intercode Environment>} \\
        \bottomrule
    \end{tabular}
    \vspace{3pt}
    \caption{Initial, Query, and Retry Prompts for Capture the Flag Evaluation.}
    \label{appx:table:prompt-ctf}
\end{table}

\begin{table}
    \centering
    \small
    \begin{tabular}{>{\raggedright\arraybackslash\tt}p{\textwidth}<{}}
        \toprulec
        \multicolumn{1}{c}{\cellcolor{lightgray} Plan Message}\\
        \midrulec
        For the following user question, let's first understand the problem and devise a plan to solve the problem. Then, let's carry out the plan to solve the problem step by step. \\
        \\
        Your plan should describe a sequence of {\color{blue} \{self.language\}} queries you can write to determine the answer. Here are three examples of coming up with a plan for a question. \\
        \\
        Question: What are the names and grades for each high schooler? \\
        Plan: \\
        1. Check what tables are available for use. \\
        2. Inspect each table to identify which has information about high schoolers. \\
        3. Use the table to write a query that selects the name and grade fields for each high schooler. \\
        ...
        \vspace{0.5em} \\\toprulec
        \multicolumn{1}{c}{\cellcolor{lightgray} Execute Plan Message}\\
        \midrulec
        You will now execute your own plan. Interact with a {\color{blue} \{self.setting\}} system using {\color{blue} \{self.language\}} queries to answer a question. Per turn, you will be given the following information: \\
        \\
        \textasciigrave\textasciigrave\textasciigrave \\
        Observation: Standard output from executing previous instruction \\
        Step: Current step \\
        \textasciigrave\textasciigrave\textasciigrave \\
        \\
        Your response should be {\color{blue} \{self.language\}} code, nothing else, formatted as follows: \\
        \textasciigrave\textasciigrave\textasciigrave{\color{blue} \{self.language\}} \\
        Your {\color{blue} \{self.language\}} code here \\
        \textasciigrave\textasciigrave\textasciigrave \\\toprulec
        \multicolumn{1}{c}{\cellcolor{lightgray} Observation Message}\\
        \midrulec
        Output: {\color{blue}<Observation from Intercode Environment>} \\
        Step: {\color{blue}<Next step to execute from the plan>}
        \vspace{0.5em} \\\toprulec
        \multicolumn{1}{c}{\cellcolor{lightgray} Post-Plan Refinement Message}\\
        \midrulec
        You have finished executing the plan, but it seems like there are still issues with your answer. Please continue to work on getting the correct answer. Per turn, you will be given the following information: \\
        \\
        \textasciigrave\textasciigrave\textasciigrave \\
        Observation: Standard output from executing previous instruction \\
        \textasciigrave\textasciigrave\textasciigrave \\
        \\
        Your response should be {\color{blue} \{self.language\}} code, nothing else, formatted as follows: \\
        \textasciigrave\textasciigrave\textasciigrave{\color{blue} \{self.language\}} \\
        Your {\color{blue} \{self.language\}} code here \\
        \textasciigrave\textasciigrave\textasciigrave \\
        \bottomrule
    \end{tabular}
    \vspace{3pt}
    \caption{Initial, Query, and Retry Prompts for Plan \& Solve Strategy Evaluation. Additional Plan \& Solve style demonstrations written for SQL and Bash can be found on the \benchmark{} project GitHub. Note that the Post-Plan Refinement Message is only used for the Plan \& Solve + Refine strategy discussed in \S~\ref{appx:experiments:additional_strategies}. It is not used for the original Plan \& Solve strategy.}
    \label{appx:table:prompt-plan-solve}
\end{table}

\end{document}
